\documentclass[11pt,a4paper,twoside]{tesis}
\usepackage[T1]{fontenc} % Use 8-bit encoding that has 256 glyphs
\usepackage[utf8]{inputenc}
\usepackage{graphicx}
\usepackage[utf8]{inputenc}
\usepackage[spanish]{babel}
\usepackage[left=3cm,right=3cm,bottom=3.5cm,top=3.5cm]{geometry}

\usepackage{amssymb}
\usepackage{amsfonts,amsmath,amsthm} % Math packages
\usepackage{mathabx}
\usepackage{mathtools}
\usepackage{tikz}
\usetikzlibrary{babel,automata,positioning,arrows}
\usepackage{graphicx}
\usepackage[rightcaption]{sidecap}
\usepackage{pgfplots}
\pgfplotsset{compat=newest}
\usepackage{xcolor}
\usepackage{paralist}
\usepackage[shortlabels]{enumitem}
\usepackage{xypic}
\usepackage[lofdepth,lotdepth]{subfig}
\usepackage{float}
%\usepackage{biblatex}
% \usepackage{txfonts}
% \usepackage{todonotes}
\newcommand*{\corner}{\mbox{\LARGE{$\mathrlap{\cdot}\lrcorner$}}}
% \usepackage{sectsty} % Allows customizing section commands
% \allsectionsfont{\centering \normalfont\scshape} % Make all sections centered, the default font and small caps

\usepackage{xspace}
\usepackage{xcolor}
\usepackage{todonotes}


\usepackage{fancyhdr} % Custom headers and footers
\pagestyle{fancyplain} % Makes all pages in the document conform to the custom headers and footers
\fancyhead{} % No page header - if you want one, create it in the same way as the footers below
\fancyfoot[L]{} % Empty left footer
\fancyfoot[C]{} % Empty center footer
\fancyfoot[R]{\thepage} % Page numbering for right footer
\renewcommand{\headrulewidth}{0pt} % Remove header underlines
\renewcommand{\footrulewidth}{0pt} % Remove footer underlines
\setlength{\headheight}{13.6pt} % Customize the height of the header

\numberwithin{equation}{section} % Number equations within sections (i.e. 1.1, 1.2, 2.1, 2.2 instead of 1, 2, 3, 4)
\numberwithin{figure}{section} % Number figures within sections (i.e. 1.1, 1.2, 2.1, 2.2 instead of 1, 2, 3, 4)
\numberwithin{table}{section} % Number tables within sections (i.e. 1.1, 1.2, 2.1, 2.2 instead of 1, 2, 3, 4)

\setlength\parindent{0pt} % Removes all indentation from paragraphs - comment this line for an assignment with lots of text

\newcommand{\osred}{\rightarrow}                % one-step reduction relation
\newcommand{\TRANS}[1]{\xrightarrow{#1}}
\newcommand{\R}{\osred}
\newcommand{\msred}{\osred^\ast}     % multi-step reduction
\newcommand{\RR}{\msred}
\renewcommand{\vec}[1]{\mathbf{#1}}
\newcommand{\SGV}{N}
\newcommand{\RSN}[2]{\RS_#1(#2)}
\newcommand{\RSB}[1]{\RSN{1}{#1}}
\newcommand{\TSdelta}{\hat{\delta}}

\newcommand{\sgs}{n}
\newcommand{\evset}{\mathcal{E}}
\newcommand{\evtequiT}[1]{\evtequi}
\newcommand{\evtequi}{\bowtie}
\newcommand{\trsysL}[1]{\xrightarrow{#1}}
\newcommand{\qst}{\;\colon\;} %such that
\newcommand{\STSE}{\mathit{TS}} % Synchronous TS abbrev
\newcommand{\SG}[1]{\STSE(#1)} % Synchronous TS (retro compa)
\newcommand{\trstabL}[1]{\stackrel{#1}{\Rightarrow}}
\newcommand{\trstab}{\trstabL{}}
\newcommand{\tra}{\pi}




\newcommand{\Q}{\mathbf{\mathsf{Q}}}
\newcommand{\C}{\mathbf{\mathsf{C}}}
\newcommand{\msgs}{\mbox{\textcolor{blue}{$\mathsf{Msg}$}}}
\newcommand{\amsg}[1][m]{\textcolor{blue}{\mathsf{#1}}}
\newcommand{\aport}[1][{}]{\textcolor{blue}{\pi^{#1}}}

\newcommand{\grmeq}{\; ::= \;}
\newcommand{\grmor}{\; \big| \;}
\newcommand{\const}[1]{\textcolor{orange}{\mathsf{#1}}}
\newcommand{\mykeyword}[1]{\textsf{\textcolor{orange}{\upshape \textbf{#1}}}}
\newcommand{\fvar}[1]{\mathrm{fv}\big(#1\big)}
\newcommand{\truek}{\mykeyword{true}}
\newcommand{\falsek}{\mykeyword{false}}
\newcommand{\ENTAILS}{\supset}
\newcommand{\iv}[1][x]{\textcolor{orange}{#1}}
\newcommand{\bool}{\textsf{Bool}}
\newcommand{\intTp}{\mathsf{Int}}



% ------------
% theorems

%\newtheorem{example}{Example}
%\newtheorem{definition}{Definition}
\newtheorem{observation}{Note}
% end theorems
% ------------

\newcommand{\actset}{A}

\newcommand{\arn}{\ensuremath{\sf ARN}\xspace}
\newcommand{\arnint}[1]{\ensuremath{\mathcal{I}_{#1}}}

\newcommand{\anarn}{\ensuremath{\alpha =  \left(X, P, C, \gamma, M_X, M_C, \mu, \Lambda\right)}\xspace}
\newcommand{\Sit}{\textcolor{blue}{\ensuremath{\mathit{S}}}}
\newcommand{\vectorize}[1]{\ensuremath{\overset{\rightarrow}{\mathbf{#1}}}}
\newcommand{\sysconf}[2]{\ensuremath{(\vectorize{#1}, \vectorize{#2})}}

%\newcommand{\cfsmalphabet}{\ensuremath{\mathcal{A}}}
\newcommand{\cfsmalphabet}{\msgs}
\newcommand{\cfsmfinitewords}{\ensuremath{\cfsmalphabet^\star}}

\newcommand{\machine}[1]{\ensuremath{\mathsf{#1}}\xspace}
\newcommand{\machp}{\machine{p}}
\newcommand{\machq}{\machine{q}}
\newcommand{\machs}{\machine{s}}
\newcommand{\machr}{\machine{r}}
\newcommand{\machx}{\machine{x}}
\newcommand{\machc}{\machine{c}}
\newcommand{\macha}{\machine{a}}

\newcommand{\FS}{\ensuremath{\mathcal{F}}\xspace}
\newcommand{\aut}[1]{\ensuremath{\textgoth{#1}}\xspace}

\newcommand{\reduct}[2]{\ensuremath{#1\hspace{-0.4em}\downharpoonccw_{#2}}\xspace}

\newcommand{\Participants}{\mathcal{P}}

\newcommand{\PSet}{\textcolor{orange}{\mathsf{P}}}
\newcommand{\chanset}{C}
\newcommand{\ptp}[1]{\mathsf{#1}}
\newcommand{\p}{\ptp p}
\newcommand{\q}{\ptp q}
\newcommand{\s}{\ptp s}
\newcommand{\rr}{\ptp r}
\newcommand{\ASigma}{\Sigma}
\newcommand{\action}{\ell}
\newcommand{\silent}{\tau}
\newcommand{\chan}[1]{\ptp{#1}}
\newcommand{\PSEND}[2]{\chan{#1} \oact \msgsort{#2}}
\newcommand{\PRECEIVE}[2]{\chan{#1} \iact \msgsort{#2}}
\newcommand{\oact}{\ptp !}
\newcommand{\iact}{\ptp ?}
\newcommand{\msgsort}[1]{\mathsf{#1}}
\newcommand{\Act}{\mathit{Act}}
\newcommand{\Inter}{\mathit{Int}}

\newcommand{\Async}[1]{\mathcal{A}(#1)}
\newcommand{\Sync}[1]{\mathcal{S}(#1)}

\newcommand{\tset}{\to}
\newcommand{\tguard}{\mathcal{G}}
\newcommand{\GSet}{\mathbb{G}}
\newcommand{\Nats}{\mathbb{N}}
\newcommand{\VSet}{\mathcal{V}}
\newcommand{\incfunc}{\mathcal{I}}
\newcommand{\Kfunc}{\mathcal{K}}

\newcommand{\trans}[2][{}]{\,\xrightarrow{#2}_{#1}\,}
\newcommand{\csconf}[2]{\conf{\vec{#1} \ ; \ \vec{#2}}}

\newcommand{\gcsconf}[4]{\conf{\vec{#1} \ ;\  \vec{#2} \ ; \ \vec{#3} \ ; \ \mathcal{#4}}}

\newcommand{\cscfg}[2]{\conf{{#1} \ ; \ {#2}}}
\newcommand{\gcscfg}[3]{\conf{{#1} \ ; \ {#2}\ ; \ {#3}}}

\newcommand{\ith}[2]{\vec{#1}[\ptp{#2}]}
\newcommand{\conf}[1]{\langle #1 \rangle}
\newcommand{\TRANSS}[1]{{\xRightarrow{\raisebox{-.3ex}[0pt][0pt]{\scriptsize $#1$} }}}
\newcommand{\Smv}{\vec{S}}
\newcommand{\Conf}[1]{\mathsf{C}_{#1}}

\newcommand{\instantiate}{\mathit{inst}}

\newcommand{\NUL}{\varepsilon}
\newcommand{\RS}{\mathbf{RS}}
\newcommand{\st}{\ \big| \ }
\newcommand{\defref}[2][{}]{Definition{#1}~\ref{#2}}
\newcommand{\secref}[2][{}]{Section{#1}~\ref{#2}}
\newcommand{\gnodes}{V}
\newcommand{\gedges}{A}
\newcommand{\glabel}{\Lambda}
\newcommand{\LSet}{\mathsf{L}}
\newcommand{\mmdef}{:=}
\newcommand{\asend}[3]{\ptp{#1} \rightarrowtriangle\ptp{#2} : \ptp{#3}}
\newcommand{\spipe}{\; | \;}

\newcommand{\ccsprc}[1]{\mathsf{#1}}
\newcommand{\Cell}{\ccsprc{Cell}}
\newcommand{\Buffer}{\ccsprc{Buffer}}
\newcommand{\Ping}{\ccsprc{Ping}}
\newcommand{\Pong}{\ccsprc{Pong}}
\newcommand{\ccsin}[2]{\mathsf{#1}?{\mathsf{#2}}}
\newcommand{\ccsout}[2]{{\mathsf{#1}!{\mathsf{#2}}}}

\newcommand{\efb}[2]{\mathcal{E}(#1,#2)}

\newcommand{\chosyn}{\fxfatal{chosyn command is not defined }}

\newcommand{\context}{\textit{context}\xspace}
\newcommand{\powerset}[1]{\ensuremath{2^{#1}}\xspace}

\renewcommand*\rmdefault{cmr}

\usepackage{xspace}
\usepackage{xcolor}
\usepackage{color}
\usepackage{todonotes}
\definecolor{LavanderPosta}{HTML}{ebd9fc}
\newcommand{\enote}[1]{\todo[fancyline,color=LavanderPosta]{#1 \\{\footnotesize(Eze)}}\xspace{}}
\newcommand{\cnote}[1]{\todo[fancyline,color=LavanderPosta]{#1 \\{\footnotesize(Charlie)}}\xspace{}}


\def \({\left (}
\def \){\right )}
\def \[{\left [}
\def \]{\right ]}
\def \[[{\left \llbracket}
\def \]]{\right \rrbracket}
\def \<{\left\langle}
\def \>{\right\rangle}

\newcommand\HeteroGenius{{{\ffmfamily Hetero}{\LobsterTwo Genius}}}
\newcommand\SEArch{{\ECFAugie SEArch}}
\newcommand\SEArchComp{{\ECFAugie SEArch -- Service Execution Architecture}}

\newcommand{\defname}[1]{\textbf{\emph{[#1]}}}
\newcommand{\theoname}[1]{\textbf{\emph{[#1]}}}
\newcommand{\propname}[1]{\textbf{\emph{[#1]}}}
\newcommand{\romboeqb}{\underline{\Diamond}}
\newcommand{\romboeqn}{\underline{\blacklozenge}}

\pgfkeys{/pgfplots/Axis Style/.style={
width=7.5cm, height=8cm,
axis x line=center, 
axis y line=middle, 
samples=100,
ymin=-1.5, ymax=5.5,
xmin=-7.0, xmax=7.0,
domain=-2*pi:2*pi,
title style={at={(1,0.5)}}}}
\newtheorem{theorem}{Teorema}
\newtheorem{lemma}{Lema}
\newtheorem{ejemplo}{Ejemplo}

\usepackage{bbm}
\def \nat{\mathbbm{N}}

\begin{document}
\def\autor{Ezequiel Davidovich Caballero}
\def\tituloTesis{Automata for Partial Binding of Services}
\def\runtitulo{Autómatas para el binding parcial de servicios}
\def\runtitle{Automata for Partial Binding of Services}
\def\director{Carlos G. Lopez Pombo}
\def\codirector{Ignacio Vissani}
\def\lugar{Buenos Aires, 2020}
\newcommand{\HRule}{\rule{\linewidth}{0.2mm}}
%
\thispagestyle{empty}

\begin{center}\leavevmode

\vspace{-2cm}

\begin{tabular}{l}
\includegraphics[width=2.6cm]{logofcen.pdf}
\end{tabular}


{\large \sc Universidad de Buenos Aires

Facultad de Ciencias Exactas y Naturales

Departamento de Computaci\'on}

\vspace{6.0cm}

%\vspace{3.0cm}
%{
%\Large \color{red}
%\begin{tabular}{|p{2cm}cp{2cm}|}
%\hline
%& Pre-Final Version: \today &\\
%\hline
%\end{tabular}
%}
%\vspace{2.5cm}

\begin{huge}
\textbf{\tituloTesis}
\end{huge}

\vspace{2cm}

{\large Tesis de Licenciatura en Ciencias de la Computaci\'on}

\vspace{2cm}

{\Large \autor}

\end{center}

\vfill

{\large

{Director: \director}

\vspace{.2cm}

{Codirector: \codirector}

\vspace{.2cm}

\lugar
}

\newpage\thispagestyle{empty}

\frontmatter
\pagestyle{empty}
%\begin{center}
%\large \bf \runtitulo
%\end{center}
%\vspace{1cm}
\chapter*{\runtitulo}

Los sistemas distribuidos resultantes de paradigmas emergentes como \emph{service-oriented computing}(SOC) y Cloud/Fog computing están transformando al mundo del software, impulsando lo que se denomina economía de API. La idea subyacente de economía de API es la posibilidad de construir software componiendo servicios externos y registrados previamente en repositorios. Aplicaciones que corren sobre recursos disponibles a nivel global con una infraestructura de comunicación que se reconfiguran en forma dinámica y transparente, a través de un middleware dedicado capaz de descubrir y conectarlas a servicios que puedan cumplir determinados requisitos. 
En general los aspectos principales del comportamiento de una API son documentados informalmente limitando la posibilidad de lograr la utopía de SOC: la negociación de servicios en forma automática. Un elemento clave para esto es la existencia de lenguajes formales, junto con técnicas asociadas de análisis, capaces de expresar por completo el contrato de comportamiento de una API. Estos formalismos usualmente están definidos de forma que la correctitud se reduce a la ausencia de ciertas configuraciones (\emph{deadlock}, \emph{receptor no especificado}, \emph{mensajes huérfanos}), que sólo puede ser comprobada con la presencia de todos los participantes. Entonces a la hora de formalizar el proceso de asegurar la interoperabilidad entre más de dos require points, todos los componentes deben comprometerse a conectarse aún cuando no se requiera su participación al principio de la comunicación. Generalized Multi-party Compatibility (GMC) es una condición suficientepara la condición de correctitud arriba detallada. En este trabajo estudiamos: 
\begin{inparaenum}[1)] 
\item una nueva clase de CFSM, llamada Multichannel Communicating Finite State Machines (mCFSMs), con la posibilidad de tener más de un canal de comunicación con los otros participantes, 
\item una definición de la propiedad de GMC para sistemas de mCFSMs,
\item una clase de Autómatas de Comunicación Asincrónica (ACAs) con la capacidad de internalizar la comuniación como operaciones de lectura/escritura en canales internos, ermitiendo la composición parcial de autómatas de comunicación y
\item un método para mapear un ACA a una mCFSM con un mecanismo para comprobar la propiedad de GMC para la clase de ACA.
\end{inparaenum}

\bigskip

\noindent\textbf{Palabras claves:} Autómatas, CFSM, SOC, GMC formalismos.

\cleardoublepage
%\begin{center}
%\large \bf \runtitle
%\end{center}
%\vspace{1cm}
\chapter*{\runtitle}

Distributed software resulting from emerging paradigms such as \emph{service-oriented computing} (SOC) and Cloud/Fog computing are transforming the world of software systems, giving impulse to what is called the API's economy. The underlying idea of it is that it is possible to construct software artifacts by composing services provided by third parties and previously registered in repositories. Applications running over globally available computational resources and communication infrastructure dynamically and transparently reconfigured, at run-time, by the intervention of a dedicated middleware capable of discovering and binding a running application with a certain requirement, to a service capable of fulfilling it.

In general the most important aspects of an API's behaviour are documented informally limiting the possibility of achieving SOC's utopia: automatic brokering of services. A key element in this quest is the existence of formal languages, together with associated analysis techniques, capable of fully expressing the API behavioral contract. These formalisms are usually defined in a way that correctness of the communication usually reduced to the absence of certain configurations (\emph{deadlock}, \emph{unspecified reception}, and \emph{orphan message}) can only be asserted in the presence of all the participants involved via properties like \emph{Generalized Multiparty Ccompatibility} (GMC).

In this work we study:
\begin{inparaenum}[1)]
 \item a new class of CFSMs, called \emph{Multichannel Communicating Finite State Machines -- mCFSMs}, with an explicit definition of the communication channels enabling, for a participant, the possibility of having more than one channel with the other participants
 \item a definition of the GMC property for systems of mCFSMs, 
 \item a class of \emph{Asynchronous Finite Communicating Automatas} (AFCAs) with the capability of internalising the communication as read / write operations on internal buffers, enabling partial composition of communicating automata, and
 \item a method for mapping an AFCA to an mCFSM providing a checking mechanism of the GMC property for the class of AFCA.
 \end{inparaenum}


\bigskip

\noindent\textbf{Keywords:} Automata, CFSM, SOC, GMC, formalisms. % OPCIONAL: comentar si no se quiere

\cleardoublepage
\chapter*{Agradecimientos}

\cnote{Escribir.} \cnote{acá van los agradecimientos} % OPCIONAL: comentar si no se quiere

\cleardoublepage
\cnote{Escribir.}\hfill \textit{A mi persona favorita.}
  % OPCIONAL: comentar si no se quiere

\cleardoublepage
\tableofcontents
\mainmatter
\pagestyle{headings}

%!TEX root = ./main.tex
\section{Introducción}
La teoría de autómatas es el estudio de artefactos abstractos de computación o "máquinas". Antes de la existencia de computadoras Alan Turing estudió una máquina abstracta que tenía todas las capacidades que las computadoras de hoy en el aspecto de qué podían computar. El objetivo de Turing era describir con precisión el límite entre lo que una computadora podía hacer y lo que no. Estas conclusiones no aplican a sus máquinas de Turing abstractas sino a las máquinas reales de hoy en día.

Entre los 40s y los 50s se estudiaron máquinas más sencillas, que hoy denominamos "autómatas finitos". Estos autómatas, originalmente pensados para modelar el funcionamiento del cebebro, resultaron ser muy útiles para una variedad de propósitos. Noam Chomsky inició el estudio de gramáticas formales. Si bien no eran estrictamente máquinas, estas gramáticas tienen una realción muy cercana con los autómatas abstractos y hoy en día sirven como la base de algunos componentes de software importantes, incluyendo partes de compiladores.

Todos estos desarrollos teóricos tienen relación directa con lo que se estudia en ciencias de la computación. Algunos de estos conceptos, como autómatas finitos y ciertas gramáticas formales, se usan en el diseño y la construcción de distintos tipos de software.

Los autómatas consisten de estados (representados gráficamente con círculos) y transiciones (representadas con flechas). Cuando el autómata ve el símbolo de entrada, ejecuta la transición a otro estado, de acuerdo con su función de transición que toma el estado actual y el símbolo reciente como parámetros de entrada. Por lo general cuando nos referimos a autómatas finitos, hablamos del conjunto de estados distintos que lo componen. Un autómata finito puede tener uno o más estados finales o accepting states. Estos estados representan el fin de la ejecución del proceso modelado por el autómata.

Los aútomatas finitos son un modelo útil para una gran variedad de software y hardware. Entre las más importantes encontramos: software para diseñar y chequear el comportamiento de circuitos digitales; el analizador lexicográfico de un compilador; software para leer textos y encontrar patrones; y, lo más relevante a este trabajo, software para verificar sistemas de todo tipo que tienen un número finito de estados distintos, como protocolos de comunicación.

La tendencia hacia los sistemas distribuidos y las redes a medida que va creciendo genera la necesidad de protocolos de comunicación más complejos. Para manejar esta complejidad se han ido introduciendo métodos formales de especificación y análisis. Dentro de estos métodos, los autómatas finitos resultan una familia de modelos apta para representar el funcionamiento de ciertos sistemas distribuidos.

Service oriented computing (SOC) es un paradigma interdisciplinario para computación distribuida. SOC cambió el modo en que las aplicaciones de software se diseñan, envían y consumen. El corazón del paradigma son servicios que proveen elementos computacionales autónomos, independientes de plataforma. Estos pueden ser descritos, publicados, descubiertos y programados usando protocolos estándar para construir redes de aplicaciones que colaboran entre sí, distribuidas dentro de distintos fronteras organizacionales.

% Service-oriented computing is an emerging cross-disciplinary paradigm for distributed computing, which is changing the way software applications are designed, delivered and consumed. At the heart of service-oriented computing are services that provide autonomous, platform-independent, computational elements that can be described, published, discovered, orchestrated and programmed using standard protocols to build networks of collaborating applications distributed within and across organizational boundaries.
 
%!TEX root = ./main.tex
\section{Preliminares}

\label{preliminares}
A continuación presentaremos algunas definiciones y resultados preliminares que serán de utilidad en las restantes secciones de esta tesis.

\begin{definition}[Clausura reflexo-transitiva] Sea $\rightarrow$ una relación sobre un conjunto S. La clausura reflexo-transitiva de $\rightarrow$, denotada  $\rightarrow^*$, se define como la relación reflexiva y transitiva más pequeña sobre S que contiene a $\rightarrow$. 
\end{definition} 
\begin{definition}[Relación de equivalencia] Es una relación binaria reflexiva, simétrica y transitiva. La relación ``es igual a''  es el ejemplo canónico donde para todo trío de objetos a, b y c, vale: 

\begin{enumerate}
    \item a= a (propiedad reflexiva)
    \item Si a = b entonces b = a (propiedad de simetría)
    \item Si a = b y b = c entonces a = c (propiedad transitiva)
\end{enumerate}
\end{definition}

\begin{prop}
 La clausura reflexo-transitiva de una relación binaria es una relación de equivalencia
\end{prop}

\subsection{Comunicación Asincrónica}
Como se dijo en la primera sección el paradigma de SOC plantea un sistema distribuido donde la comunicación entre distintos componentes se realiza en forma asincrónica. El sistema en tiempo de ejecución se conecta con los distintos servicios a medida que los va requiriendo. El intercambio de mensajes con cada componente es independiente entre sí, salvo que exista alguna necesidad de lo contrario. Para representar este comportamiento utilizamos un tipo de autómata finito denominado \emph{Communicating Finite State Machines} (llamado a partir de ahora CFSM). El concepto de CFSM fue introducido en \cite{CFSM} con el objetivo modelar y estudiar el comportamiento de sistemas distribuidos constituídos por un conjunto de procesos secuenciales que ejecutan concurrentemente y se comunican a partir de intercambiar mensajes a través de canales de comunicación previamente declarados. Con cada CFSM representamos los distintos componentes del sistema distribuido. 

La naturaleza dinámica no pre programada
\begin{definition}[Communicating Finite State Machines] Sea $\mathcal{M}$ un conjunto finito de mensajes y un conjunto finito de mensajes  y un conjunto finito de participantes $\mathcal{P}$ un conjunto finito de participantes, definimos una CFSM sobre $\mathcal{M}$ como un sistema de transición finito $\langle Q, C, q_0, \mathcal{M}, \delta \rangle$ donde
\begin{itemize}
  \item $Q$ es un conjunto finito de estados;
  \item $C = \left\{ pq \in \mathcal{P}^2 \left|\right. p \not= q\right\}$ es un conjunto de canales
  \item $q_0 \in Q$ es el estado inicial;
  \item $\delta \subseteq Q \times (C \times \{!,?\} \times \mathcal{M}) \times Q$ es un conjunto finito de \emph{transiciones}.
  \end{itemize}
  
\end{definition} 

A partir de la introducción de las CFSMs, la siguiente definición introduce el concepto de \emph{Communicating System} \cite[Def.2.2]{Communicating System} como un conjunto de CFSMs que cumplen determinadas propiedades.

\begin{definition}[Communicating System] Dada una CFSM $\textit{M}_p = \langle Q_p, q_{0_p}, \mathcal{M}, \delta_p \rangle $ para cada participante $p \in \mathcal{P}$, la tupla $S= \langle M_p \rangle_{p \in \mathcal{P}} $ es un communicating system CS.
Una configuración de S es un par $s = \langle \overrightarrow{q} ; \overrightarrow{w} \rangle$ donde $\overrightarrow{q} = (q_p)_{p \in \mathcal{P}}$ con $q_p \in Q_p$ y donde $\overrightarrow{w} = (w_{pq})_{pq \in C}$ con $ w_{pq} \in \mathcal{M}^*$. La componente $\overrightarrow{q}$ es el estado de control y $q_p \in Q_p$ es el estado local de la máquina $ M_p$. La configuración inicial de S es $s_0 = \langle \overrightarrow{q_0} ; \overrightarrow{\epsilon} \rangle$ con $\overrightarrow{q_0} = (q_{0_p})_{p \in \mathcal{P}}$ y $\overrightarrow{\epsilon} = (\epsilon_p)_{p \in \mathcal{P}}$
% Un \emph{communicating system} es un mapa \textit{S} que asigna una CFSM $S(p)$ a cada $p \in \mathcal{P}$. Escribimos $q \in S(p)$ cuando $q$ es un estado de la máquina $S(p)$ y $t \in S(p)$ cuando $t$ es una transición de $S(p)$.\\

\end{definition}

La semántica de los CSs está dada por un sistema de transición etiquetado cuyos estados y transiciones determinan las posibles ejecuciones del conjunto de procesos sobre el que está definido el sistema.


\begin{definition}[Estados y configuraciones alcanzables]
\label{def:estadosyconf}
  Una configuración $c' = \langle \overrightarrow{q'} ; \overrightarrow{w'} \rangle$ es \emph{alcanzable} desde otra
  configuración $c = \langle \overrightarrow{q} ; \overrightarrow{w} \rangle$ a través de la \emph{ejecución de la transición} $l$ (escrito $s \overset{l}{\rightarrow} s'$) si existe un $m \in \mathcal{M}$ tal que ocurre una de las siguientes alternativas:
	\begin{enumerate}
		\item $l = sr!m$, $\langle q_s, l,  q'_s\rangle \in \delta_s$ y 
			\begin{enumerate}
				\item $q'_p = q_p$ para todo $\p \neq s$; y
				\item $w'_{sr} = w_{sr} \cdot m$ y  $w'_{pq} = w_{pq}$ para todo $pq \neq sr$; o bien
			\end{enumerate}
		\item $l = sr?m$, $\langle q_r, l,  q'_r\rangle \in \delta_r$ y 
			\begin{enumerate}
			\item $q'_{p} = q_{p}$ para todo $p \neq r$; y
				\item $w_{sr} = m \cdot w'_{pq}$ y $w'_{pq} = w_{pq}$ para todo $pq' \neq sr$
			\end{enumerate}
	\end{enumerate}
Escribimos $ s_1 \overset{t1...tm}{\rightarrow} s_{m+1}$ cuando para algún $s_2,...,s_m, s_1\overset{t}{\rightarrow} s_2...s_m\overset{t_m}{\rightarrow} s_{m+1} $ . El conjunto de configuraciones alcanzables de S es $RS(S) = \{s | s_0 \rightarrow^*S \}$
\end{definition}

En adelante presentaremos definiciones y resultados que permiten expresar propiedades de los CSs introducidos en la definición anterior.

\begin{definition}[Communicating System seguro] Sea $S$ un CS; se dice que $S$ es \emph{seguro} si para cada $s \in RS(S)$:
\begin{enumerate}
\item $s$ no es una configuración de deadlock (cuando la única transición de salida necesita consumir un mensaje que no está en el buffer)
\item $s$ no posee mensajes huérfanos (cuando no hay transiciones salientes y quedaron mensajes en el buffer) , y
\item $s$ no posee recepciones no especificadas (cuanndo la única transición saliente consume un mensaje disinto al primero de la cola).
\end{enumerate}
Para poder expresar esta condición de seguridad es necesario identificar conjuntos de acciones que pueden ser llevadas a cabo concurrentemente. Para esto definimos las siguientes relaciones sobre el conjunto de transiciones de una CFSM. Dados $q, q' \in Q$, se define $\mathit{act}(q,q') = \left\{\ell \ \left|\right. \ (q,\ell,q') \in \delta \right\}$ y $\Diamond, \blacklozenge \subseteq \delta \times \delta$ como las relaciones de equivalencia más pequeñas que contienen $\romboeqb$ y $\underline{\blacklozenge}$ donde:
\begin{itemize}
\item $(q_1, \ell, q_2) \underline{\Diamond} (q'_1, \ell, q'_2)$ sii $ l \notin \mathit{act}(q_1, q'_1) \land \mathit{act}(q_1, q'_1) = \mathit{act}(q_2, q'_2) \land \mathit{act}(q_2, q'_2) \neq \emptyset $
\item  $(q_1, \ell, q_2) \underline{\blacklozenge} (q'_1, \ell, q'_2)$ sii $ (q_1, \ell, q_2) \romboeqb (q'_1, \ell, q'_2) $ y $\forall(q,\ell,q') \in [(q_1, \ell, q_2) ]^{\Diamond}, \ \mathit{act}(q_1,q) = \mathit{act}(q_2,q') \land \mathit{act}(q'_1,q) = \mathit{act}(q'_2,q')$  
\end{itemize}
donde $[(q_1, \ell, q_2) ]^{\Diamond}$ es la clase de equivalencia de $(q, \ell, q')$ respecto de la relación $\Diamond$ (resp. $\blacklozenge$). Intuitivamente dos transiciones están $\blacklozenge$-relacionadas si se refieren a la misma acción aún teniendo en cuenta el interleaving.
\end{definition}

\begin{example}[Relaciones $\underline{\Diamond}$ y $\underline{\blacklozenge}$]
\label{ex:relaciones}
Consideremos la siguiente CFSM:
\begin{center}
\begin{tikzpicture}[->, thick]
 \node[state,initial] (q_0)   {$q_0$}; 
 \node[state] (q_1) [right= of q_0 ] {$q_1$};
 \node[state] (q_5) [right= of q_1 ] {$q_5$};
 \node[state] (q_2) [below= of q_0 ] {$q_2$};
 \node[state] (q_3) [right= of q_2 ] {$q_3$};
 \node[state,accepting] (q_6) [right= of q_3 ] {$q_6$};
 \draw[]        
        (q_0) edge[above] node{sr!a} (q_1)
        (q_0) edge[right] node{sr'!b} (q_2)
        (q_1) edge[right] node{sr'!b} (q_3)
        (q_1) edge[above] node{sr!a} (q_5)
        (q_2) edge[above] node{sr!a} (q_3)
        (q_2) edge[bend right, below] node{sr!c} (q_6)
        (q_3) edge[above] node{sr!a} (q_6)
        (q_5) edge[right] node{sr'!b} (q_6)
        ;
\end{tikzpicture} 
\end{center}

\begin{enumerate}
    \item \label{ex:cond1} $(q_0, sr!a,q1) \underline{\Diamond} (q_2,sr!a,q3)  $ %dado que está entrelazada con $sr'!b$
    \item \label{ex:cond2} $(q_0, sr!a,q1) \underline{\blacklozenge} (q_2,sr!a,q3) $ %por la misma razón que 1
    \item \label{ex:cond3} No vale $ ((q_0, sr!a,q1) \underline{\Diamond} (q_1,sr'!b,q5))$ %porque la transición entre $q_0$ y $q_1$ pasa a través de $sr!a$. Es decir que las dos transiciones son secuenciales no concurrentes. 
    \item \label{ex:cond4} $ (q_0, sr'!b,q2) \underline{\Diamond} (q_1, sr'!b,q3)$ %la relación e 
    \item \label{ex:cond5} No vale $((q_0, sr'!b,q2) \underline{\blacklozenge} (q_1,sr'!b,q_3)) $
\end{enumerate}

Las relaciones en Ej.~\ref{ex:relaciones}.\ref{ex:cond1}--~\ref{ex:relaciones}.\ref{ex:cond2} se sostienen dado que ambas transiciones están entrelazadas con $sr'!b$. La relacion en Ej.~\ref{ex:relaciones}.\ref{ex:cond3} no se sostiene debido a que la transición entre el origen de una $(q_0)$ y el del otro $(q_1)$ pasa por $sr!a$. Ambas transiciones en Ej.~\ref{ex:relaciones}.\ref{ex:cond3} son secuenciales, no concurrentes. La relación en Ej.~\ref{ex:relaciones}.\ref{ex:cond4} se sostiene, pero en Ej.~\ref{ex:relaciones}.\ref{ex:cond5} no porque $(q_5,sr'!b,q_6)$ está en la clase de $\Diamond$-equivalencia de $(q_0,sr'!b,q_2)$ para la cual la condición no se sostiene (debido a la transición con la etiqueta $sr!c$).
\end{example}

\begin{definition}[Eventos] Dados un conjunto de participantes $\mathcal{P}$ y un conjunto de mensajes $\mathcal{M}$ definimos un evento $e$ como es una tupla $\langle q_s, q_r, s, r, a \rangle)$ (también escrita como $\langle q_s, q_r, s \rightarrow r:a \rangle$), tal que $s,r \in \mathcal{P}$, indicando que $s$ y $r$ pueden intercambiar el mensaje a, cuando están en el estado $q_s$ y $q_r$, respectivamente.
\end{definition}

\marginpar{REvisar} Tomando en consideración la definición anterior, para distinguir el paralelsimo a nivel máquina introducimos una relación de equivalencia sobre eventos que identifica eventos cuyas transiciones son $\blacklozenge$-equivalentes.

\begin{definition}[Equivalencia entre eventos] Definimos la equivalencia entre eventos como $\bowtie = \bowtie_s \cap \bowtie_r \subseteq \mathcal{E} \times \mathcal{E}$ donde $\mathcal{E}$ es el conjunto de eventos del sistema y se cumplen las siguientes condiciones:

\begin{itemize}
\item $(q_1, q_2, s \rightarrow r:a) \bowtie_s (q'_1, q'_2, s \rightarrow r:a) \iff$ \\ 
 $\forall (q_1, sr!a,q_3),(q'_1, sr!a,q'_3) \in \delta_s:(q_1, sr!a,q_3) \underline{\blacklozenge} (q'_1, SR!a,q'_3) $
\item $(q_1, q_2, s \rightarrow r:a) \bowtie_r (q'_1, q'_2, s \rightarrow r:a) \iff \\ 
\forall (q_2, sr?a,q_3),(q'_2, sr?a,q'_4) \in \delta_ r:(q_2, sr?a,q_4) \romboeqn (q'_2, sr?a,q'_4)$ \end{itemize}
\end{definition}
\begin{example}[Equivalencia entre eventos]
\label{ex:equiveventos}
Considere el siguiente Communicating System

\begin{tikzpicture}[->, thick]
 \node[state,initial] (q_0)   {$q_0$}; 
 \node[state] (q_1) [right= of q_0 ] {$q_1$};
  \node[state] (q_2) [below= of q_0 ] {$q_2$};
 \node[state] (q_3) [right= of q_2 ] {$q_3$};

 \draw[]        
        (q_0) edge[above] node{pr!a} (q_1)
        (q_0) edge[right] node{sp?b} (q_2)
        (q_1) edge[right] node{sp?b} (q_3)
        (q_2) edge[above] node{pr!a} (q_3)
        ; 
\end{tikzpicture} 
\qquad
\begin{tikzpicture}[->, thick]
 \node[state,initial] (q_0)   {$q_0$}; 
 \node[state] (q_1) [below= of q_0 ] {$q_1$};

 \draw[]        
        
        (q_0) edge[right] node{pr?a} (q_1)
        
        ;
\end{tikzpicture} 
\qquad
\begin{tikzpicture}[->, thick]
 \node[state,initial] (q_0)   {$q_0$}; 
 \node[state] (q_1) [below= of q_0 ] {$q_1$};

 \draw[]        
        
        (q_0) edge[right] node{sp?b} (q_1)
        
        ;
\end{tikzpicture} 

En el sistema de arriba podemos ver los siguientes eventos $(q_{0p}, q_{0r}, p \rightarrow r:a)$, $(q_{2p}, q_{0r}, p \rightarrow r:a)$, $(q_{0s}, q_{0p}, s \rightarrow p:)$ y $(q_{0s}, q_{3p}, s \rightarrow p:b)$. Queremos ver si se cumple que $(q_{0p}, q_{0r}, p \rightarrow r:a) \bowtie (q_{2p}, q_{0r}, p \rightarrow r:a)$, es decir que son eqivalentes tanto bajo $\bowtie_p$ como en $\bowtie_r$. 

Para el primero queremos ver que $(q_{0p}, pr!a, q_{1p}) \romboeqn (q_{2p}, pr!a, q_{3p})$
Para esto necesitamos que valga $(q_{0p}, pr!a, q_{1p}) \romboeqb (q_{2p}, pr!a, q_{3p})$. Esto se cumple dado que $pr!a \notin act(q_{0p},q_{2p})$ y $act(q_{0p},q_{2p}) = act(q_{1p},q_{3p}) \neq \emptyset$. Ahora tenemos que ver la clase de equivalencia $[(q_{0p}, pr!a, q_{1p})]^{\Diamond}$. La misma es el conjunto unitario $ \{(q_{2p}, pr!a, q_{3p})\}$. Por último queremos ver que  $act(q_{0p}, q_{2p}) = 
act(q_{1p}, q_{3p}) \land act(q_{2p}, q_{2p}) = act(q_{3p}, q_{3p})$. La segunda es trivial, la primera se cumple siendo el conjunto unitario $\{sp?b\}$. Con esto demostramos que $(q_{0p}, q_{0r}, p \rightarrow r:a) \bowtie_p (q_{2p}, q_{0r}, p \rightarrow r:a)$ 

Nos queda probar que vale $(q_{0p}, q_{0r}, p \rightarrow r:a) \bowtie_r (q_{2p}, q_{0r}, p \rightarrow r:a)$. Esta equivalencia es más sencilla de probar, dado que tenemos una única transición. Entonces vemos que es trivial que $(q_{0r}, pr?a, q_{1r}) \underline{\blacklozenge} (q_{0r}, pr?a, q_{1r})$. Con esto vemos que $(q_{0p}, q_{0r}, p \rightarrow r:a)$ y $(q_{2p}, q_{0r}, p \rightarrow r:a)$ son equivalentes en $\bowtie_p$ y $\bowtie_r$, por lo tanto están en $\bowtie = \bowtie_p \cap \bowtie_r$.
\end{example} 

A continuación definimos la noción de Sistema de Transición Sincrónico, para reflejar el comportamiento de un sistema cuando envíos y recepciones son emparejados para mostrar que ocurren al mismo tiempo.
\newpage

\begin{definition}[Sistema de transición sincrónico] Dados un $S = (M_P)_{p \in \mathcal{P}}$ un CS, sea $\langle N,\hat{\delta}, E \rangle$, donde: 
\begin{itemize}
    \item[] $N = \{\overrightarrow{q} \ | \ (\overrightarrow{q}; \overrightarrow{\epsilon}) \in RS_1(S) \}$,
    \item[] $\hat{\delta}= \{(n, e, n') \ | \ (n;\overrightarrow{\epsilon}) s_1 \overset{sr!a}{\longrightarrow}\overset{sr?a}{\longrightarrow} (n';\overrightarrow{\epsilon})	\land e= n[s], n[r], s \rightarrow r:a \}$, y
    \item[] $ E = \{ \exists n, n' \in N : (n,e,n') \in \hat{\delta}\} \subseteq \mathcal{E}$,
\end{itemize}
   el \emph{Sistema de Transición Sincrónico} de S es $TS(S)= \langle N, n_0, E/ \bowtie,\rightrightharpoons \rangle$ donde $n_0= \overrightarrow{q_0} $ es el estado inicial, $n \overset{[e]}{\rightrightharpoons} n' \iff (n,e,n') \in \hat{\delta}$. Fijamos un conjunto $\hat{E}$ de elementos representativos de cada clase de equivalencia $\bowtie$ (ej: $\hat{E} \subseteq E$ y $\left(\forall e \in E\right)\left(\exists!e' \in \hat{E}\right)\left(e' \in [e] \right)$) y escribimos $n \overset{e'}{\rightrightharpoons} n'$ para $ n \overset{[e]}{\rightrightharpoons} n'$ cuando $ e' \in [e] \cap \hat{E} $. Las secuencias de eventos se notan con un símbolo $\pi$ y extendemos la notación de $ \rightarrow$ en la Def.~\ref{def:estadosyconf} a $\rightrightharpoons$ (ej: $si \pi = e_1 ...e_k, n_1 \overset{\pi}{\rightrightharpoons}n_{k+1} sii n_1 \overset{e_1}{\rightrightharpoons} n_2 \overset{e_2}{\rightrightharpoons}...\overset{e_k}{\rightrightharpoons} n_{k+1}$).

$TS(S)$ representa todas las posibles ejecuciones sincrónicas del sistema $S$; y cada transición es etiquetada con un evento $e$.
\end{definition}

\begin{example}[Sistema de Transición Sincrónico] 
\label{ex:STS}
Consideremos el Ej.~\ref{ex:equiveventos}, su Sistema de Transición Sincrónico es el siguiente

\begin{tikzpicture}[->, very thick]
\node[state] (q_0) {};
\node[state] (q_1) [right= 3.5cm of q_0 ] {};
\node[state] (q_2) [below= of q_0 ] {};
\node[state] (q_3) [right= 3.5cm of q_2 ] {};
%$(q_{0p}, q_{0r}, p \rightarrow r:a)$, $(q_{2p}, q_{0r}, p \rightarrow r:a)$, $(q_{0s}, q_{0p}, s \rightarrow p:)$ y $(q_{0s}, q_{3p}, s \rightarrow p:b)$
\draw[]        
        (q_0) edge[above] node{$(q_{0p}, q_{0r}, p \rightarrow r:a)$} (q_1)++
        (q_0) edge[left] node{$(q_{0s}, q_{0p}, s \rightarrow p:b)$} (q_2)
        (q_1) edge[right] node{$(q_{0s}, q_{3p}, s \rightarrow p:b)$} (q_3)
        (q_2) edge[below] node{$(q_{2p}, q_{0r}, p \rightarrow r:a)$} (q_3)
        ; 
\end{tikzpicture}

Tenemos $(q_{0p}, q_{0r}, p \rightarrow r:a) \bowtie (q_{2p}, q_{0r}, p \rightarrow r:a)$ y $(q_{0s}, q_{0p}, s \rightarrow p:) \bowtie (q_{0s}, q_{3p}, s \rightarrow p:b)$. Podemos considerar equivalentes los eventos de las transiciones verticales por un lado y las de las transiciones horizontales por el otro. Esto nos permite identificar un par de interacciones concurrentes, pero seguir diferenciándolas de otras instancias de comunicación $p \rightarrow r:a$ y $s \rightarrow p:b$
\end{example}

\begin{definition}[Proyecciones] La proyección de un evento e sobre un participante p, denotado por $e \downharpoonright_p$ se define de la siguiente manera:

\begin{equation}
(q_s,q_r,s \rightarrow r:a) \downharpoonright_p = \begin{cases} 
pr!a & \mathit{if} s=p \\
sp?a & \mathit{if} r=p \\
\epsilon & \mathit{en \ otro \ caso} \\
\end{cases} 
\end{equation}

La proyección se define sobre secuencias de eventos en el modo evidente. La proyección $TS(S)= (N, n_0, \hat{E}, \rightrightharpoons)$ sobre el participante p, notada $ TS(S) \downharpoonright_p $, es el autómata $(Q, q_0, \Sigma, \delta)$ donde $Q=N, \ q_0 = n_0$, $\Sigma$ es el conjunto de etiquetas y $\delta \subseteq Q \times \Sigma \cup \{ \epsilon \} \times Q $ es el conjunto de transiciones, tal que $(n_1, e \downharpoonright_p, n_2) \in \delta \iff n_1 \overset{e}{\rightrightharpoons} n_2 $

\end{definition}

\begin{example}{Proyecciones}


\end{example}

Introducimos el concepto de generalized multiparty compatibility (GMC) como una condición completa y sólida para construir CFSMs. 
% Usé sólida como traducción de sound. En el paper habla de construir global graphs, término que acá no usamos, no se si puse bien
A partir de este punto, fijamos un sistema $S =(M_p)_{p \in \mathcal{P}} $ con $TS(S)= (N, n_0, \hat{E}, \rightrightharpoons)$. GMC depende de dos condiciones, representabilidad y propiedad de ramificación.\\

La propiedad de representabilidad determina que cada máquina, traza y elección estén representadas en el sistema de transición.

\begin{definition}Para un lenguaje $\mathcal{L}$, $hd(\mathcal{L})$ devuelve las primeras acciones de $\mathcal{L}$ si las tiene: 
\begin{center}
$hd(\mathcal{L})= \{\ell \ | \ \exists q \in Act*: \ell \cdot q \in \mathcal{L} \}$	\ \	 $hd(\{ \epsilon \}) = \{ \epsilon \}$
\end{center}
Dado $n \in N$, sea $ TS(S)\langle n \rangle$ el sistema de transición TS(S) donde se reemplaza al estado inicial $n_0$ por $n$. Escribimos LT(S,n,p) para $\mathcal{L}(TS(S)\langle n \rangle)\downharpoonright_p $, es decir LT(S,n,p) es el lenguaje que se obtiene estableciendo a n como nodo inicial de TS(S) y proyectando el nuevo sistema de transición sobre p.
\end{definition}

\begin{definition}[Representabilidad]Un sistema S es representable si
\begin{enumerate}
\item $\mathcal{L}(M_p) = LT(S,n_0,p) $ y
\item $\forall q \in Q_p \ \exists n \in N: n[p] = q \land  \cup_{(q,\ell, q') \in \delta_p} \{ \ell \} \subseteq hd (LT(S,n,p))$
\end{enumerate}
para todo $p \in \mathcal{P} $ \\

La primera condición asegura que cada traza de cada máquina esté en $TS(S)$, a su vez, la segunda condición es necesaria para asegurar que cada elección en cada máquina esté representada en $TS(S)$.
\end{definition}

La propiedad de ramificación estimula que cada vez que hay una decisión en un sistema $TS(S)$, una única máquina toma esa decisión y cada uno de los otros participantes es notificado de la rama que se tomó o no participa en esa elección. 
 
\begin{definition}[Propiedad de ramificación] Un sistema S posee la propiedad de ramificación si para todo $n \in N$ y para todo $e_1 \neq e_2 \in \hat{E} \ \mathit{tq} \ n \overset{e_1}{\rightrightharpoons} n_1 \ \mathit{y} \  \ n_1 \overset{e_2}{\rightrightharpoons} n_2$ luego tenemos
\begin{enumerate}
\item o bien existe $n'\in N$ tal que $n_1 \overset{e_2}{\rightrightharpoons} n' \ \mathit{y} \ n_2 \overset{e_1}{\rightrightharpoons} n'$, o 
\item para cada $(n'_1, n'_2) \in ln(n, e_1, e_2)$ quedando 

$L_p^i = hd (\{e_i \downharpoonright_p \cdot \phi \  | \ \phi \in LT(S,n'_i,p) \} ) $ con $i \in \{1,2\}$ y $p \in P$, y se cumplen las condiciones 2a, 2b y 2c definidas abajo 	

\begin{enumerate}[(a)]
\item Choice awareness: $\forall p \in \mathcal{P}$ valen \begin{enumerate}[i.]
\item $L_p^1 \cap	L_p^2 \subseteq \{ \epsilon \} \ \mathit{y} \  \epsilon \in L_p^1 \iff \epsilon \in L_p^2$, o 
\item $\exists n' \in N, \pi_1, \pi_2$: $n'_1 \overset{\pi_1}{\rightrightharpoons} n' \land  n'_2 \overset{\pi_2}{\rightrightharpoons} n' \land (e_1 \cdot \pi_1) \downharpoonright_p= (e_2 \cdot \pi_2)\downharpoonright_p = \epsilon$
\end{enumerate}
\item selector único: $\exists!s \in \mathcal{P}: L_s^1 \cap L_s^2 = \emptyset \land \exists sr!a \in L_s^1 \cup L_s^2 $
\item no race: $\forall r \in \mathcal{P}: L_r^1\cap L_r^2 = \emptyset \Rightarrow \forall s_1r?a_1 \in L_r^1, \forall s_2r?a_1 \in L_r^2:\forall i \neq j \in \{1,2\}: n'_i \overset{\pi_i}\rightrightharpoons  $ \\
$\Rightarrow dep (s_i \rightarrow r: a_i, e_i \cdot \pi_i, s_j \rightarrow r: a_j) $
\end{enumerate}
\end{enumerate} 	
\end{definition}

Representabilidad garantiza que TS(S) contiene suficiente información para decdir propiedades seguras de cualquier ejecución asincrónica de S. La propiedad de ramificación asegura que si una rama en TS(S) representa una elección esta está "bien formada".\\
 
Con esto definimos que toda ramificación es o bien (1) la ejecución concurrente de dos eventos; o, para cada participante p (2(a)I) si p no termina antes de n entonces las primeras dos acciones de p en dos ramas distintas son disjuntas; o (2(a)II) p no está involucrado en la elección, o sea la ramas se juntan antes de que p realice ninguna acción; (2(b)) hay un único participante s tomando la decisión; y (2(c)) para cada participante r involucrado en la elección, no puede haber race condition entre los mensajes que puede recibir r. La no race condition asegura que en ninguna ejecución (asincrónica) de S, si una máquina tiene más de un buffer no vacío, entonces puede leerlos en cualquier orden (interleaving es posible). Notar que si una máquina r recibe todos sus mensajes de un mismo emisor, entonces hay una $\triangleleft$-relación entre todas sus acciones.

\begin{definition}[General Multiparty Compatibilty] Un sistmea S es generalised multiparty compatible (GMC) si es representable y posee la propiedad de ramifiación. 
\end{definition}


%!TEX root = ./main.tex


\chapter{Componentes Asincrónicas}
\label{AFCA}
Como ya hemos mencionado, en SOC los sistemas son concebidos como objetos dinámicos construidos en \emph{run-time} en la medida que su ejecución llega a un estado en el que la intervención de servicios externos se hace necesaria. Es decir un sistema de este tipo utilizará distintos servicios según las necesidades que se manifiesten a lo largo de una ejecución particular. 
 
Una aplicación que se encuentra ejecutando se conecta con los servicios que le son necesarios a través de canales de comunicación por los cuales se envían o reciben mensajes. Estos canales pueden establecer una comunicación entre un número fijo (para cada canal particular) pero no acotado a priori de servicios. En la sección anterior hemos detallado un conjunto de propiedades que garantizan una comunicación sin errores a través de estos canales (i.e. ausencia de deadlock, ausencia de mensajes huérfanos y ausencia de situaciones en las que el receptor no se encuentra a la espera de un mensaje que le fue enviado) y un procedimiento para garantizarlas. Estas condiciones y su procedimiento de análisis parten de la hipótesis de que las CFSMs correspondientes a todos y cada uno de los participantes de la comunicación sobre dicho canal se encuentran disponibles.

Ahora bien, que la aplicación arribe a un estado en el que un servicio se hace necesario sobre un canal particular, no implica que todos los participantes también lo sean en ese mismo instante y por ello, con el objeto de profundizar esta concepción incremental, a demanda, que se tiene sobre los sistemas de software surge, más o menos naturalmente, la idea de poder dotar al \emph{middleware} de la capacidad de realizar un \emph{binding} parcial sobre los canales. A esta práctica la llamaremos \emph{binding incremental}.
 
Esta percepción parcial del \emph{binding} sobre un canal requiere la utilización de un un lenguaje de descripción que soporten dichos mecanismos de composición. Por ejemplo, al componer dos CFSMs puede ocurrir que cada máquina se comunique a través de un canal con un tercer participante, estos dos canales son independientes y, además, una vez que se ha realizado la composición, deben ser percibidos por este tercer participante como canales de comunicación separados. Por lo tanto, para preservar la semántica de la comunicación, es necesario que el CFSM resultante de la composición tenga dos canales con este tercer participante. Este fenómeno será el eje rector de las modificaciones que introduciremos en esta sección. Esta característica no solo será necesario a nivel de CFSM (interfaz de comunicación de un servicio) sino también de los autómatas que caracterizan el cómputo, cuya interfaz de comunicación es expresada a través de una CFSM.

\section{Autómatas finitos de comunicación asincrónica}
Las CFSMs son un tipo de autómata que modela únicamente la comunicación externa con otros participantes del sistema. Necesitábamos un tipo de autómata que además representara otros comportamientos internos. Para esto primero recurrimos a los Input Output Automata \cite{lynch:cwiq-2_3}. 
Un I/O automaton modela un componente de un sistema distribuido que puede interactuar con otros sistemas componentes. Es un tipo de máquina de estados en la cual las transiciones están asociadas con determinadas acciones específicas. Existen tres tipos de acciones: input, output, y acciones internas. El autómata usa sus acciones de input y output para comunicarse con el enotrno, mientras que las acciones internas son sólo visibles para el autómata. A diferencia de las acciones internas y las de output que son seleccionadas y resueltas por el autómata, las de entrada, que llegan del entorno, no están bajo control del autómata. Para componer estos autómatas se requiere que sincronicen las acciones compartidas. Es decir cuando un autómata ejecuta una transición con la etiqueta $\pi$, todo autómata que tenga esa acción $\pi$ la ejecuta en simultáneo.
El problema es que los IO no representan bien la comunicación asincrónica. Las operaciones de comunicación externa sincronizan a todos los participantes del sistema. Es decir que si dos autómatas componentes tienen una misma transición, se ejecutan en simultáneo. Esto va en contra de la idea de asincronsidad que buscamos.
Para modelar este comportamiento asincrónico introducimos los Autómatas Finitos de Comunicación Asincrónica (AFCA) (ver Def.~\ref{def:afca}). Estos autómatas tienen tres tipos de transiciones: 
\begin{inparaenum}[1)]
\item internas, que sirven el propósito de representar cómputo realizado por la componente; 
\item de buffer, que representan comunicación asincrónica interna entre elementos de la componente, y que permiten representar la comunicación entre dos participantes luego de una composición; y por último 
\item de entrada / salida, que modelan acciones de comunicación con otras componentes del sistema.
\end{inparaenum} \\
Como lo que queremos es modelar sistemas distribuidos con comunicación asincrónica, necesitamos una noción de composición. Para esto introducimos una definición de composición para los AFCA (ver Def.~\ref{def:composicion}). 

\begin{definition}[Cola]
\label{def:cola}
Una cola es un tipo de estructura de datos caracterizada por ser una secuencia de elementos first in-first out (FIFO) debido a que el primer elemento en entrar es el primer elemento en salir. Para manejarla se definen dos operaciones, una para agregar y otra para sacar elementos de la una cola. Para los usos de este trabajo vamos a definir colas en las cuales se depositan y de las cuales se retiran mensajes, con las operaciones:

\begin{itemize}
\item Cola vacía: denotado $[\ ]$
\item Encolar mensaje: denotado como $b \ll m$, si $b$ es una cola y $m$ es un mensaje.
\item Desencolar mensaje: denotado como $b \gg m$, si $b$ es una cola y $m$ es un mensaje.
\item Sea una cola vacía $b= [\ ]$ al aplicarle la operación $b \ll m$ queda $b= [m]$
\item Sea una cola no vacía $b= [m_1,..,m_n]$ al aplicarle la operación $b \ll m$ queda $b= [m,m_1,..,m_n]$ del mismo modo al aplicarle $b \gg m_n$ queda $b= [m,m_1,..]$
\end{itemize}
\end{definition}

\begin{definition}[Autómata finito de comunicación asincrónica]
\label{def:afca}
Sea $\mathcal{P}$ un conjunto de participantes y $\mathcal{M}$ un conjunto de mensajes. Un \emph{autómata finito de comunicación asíncrona} es una estructura $A_\mathcal{P} = \langle Q, B, \mathcal{C}, \Sigma, \delta, q_0, F\rangle$ tal que:

\begin{itemize}
\item $Q$ es un conjunto finito de estados,
\item $B \subseteq \{pq_n \ | \ pq \in \mathcal{P}^2, n \in \mathbb{N}, p \neq q \} $ es un conjunto finito de buffers (i.e. colas, ver Def.~\ref{def:cola}),
\item $\mathcal{C} \subseteq \{pq_n \ | \ pq \in \mathcal{P}^2, n \in \mathbb{N}, p \neq q \}$ es un conjunto de canales tales que $B \cap \mathcal{C} = \emptyset$ 
\item $\Sigma = \{ \Sigma_\mathit{Int} \cup \Sigma_\mathit{Ex} \cup \Sigma_\mathit{Buff}\} $, $\Sigma \cap \mathcal{M} = \emptyset$ es el conjunto de etiquetas del autómata, siendo
\begin{inparaenum}[1)]

\item $\Sigma_\mathit{Int}$ las acciones internas del autómata 

\item $\Sigma_\mathit{Ex}$ un conjunto de etiquetas de la forma $\mathit{In}(c,m), \mathit{Out}(c,m)$ dónde $c \in \mathcal{C}$ es el canal a través del cual se resuelve la misma y $m \in \mathcal{M}$ es el mensaje .

\item $\Sigma_\mathit{Buff}$ es el conjunto de etiquetas de las acciones sobre los buffers de la forma $b \ll m$ o $b \gg m$, dónde $b \in B$ y $m \in \mathcal{M}$.
\end{inparaenum}
\item $\delta = (\delta_\mathit{Int} \cup \delta_\mathit{Ex} \cup \delta_\mathit{Buff})$ siendo:

\begin{inparaenum}[1)]
\item $\delta_\mathit{Int} \subseteq Q \times \Sigma_\mathit{Int} \times Q$ es la relación de transición por acciones internas de $A_\mathcal{P}$, %transiciones internas

\item $\delta_\mathit{Ex} \subseteq Q \times \Sigma_\mathit{EX}  \times Q$ es la relación de transición de comunicación externa de $A_\mathcal{P}$, %comunicación externa

\item $\delta_\mathit{Buff} \subseteq Q \times \Sigma_\mathit{Buff} \times Q$  es la relación de transición de comunicación interna de $A_\mathcal{P}$ %comunicación interna
\end{inparaenum}
\item $q_0 \in Q$ es el estado inicial, y
\item $F \subseteq Q$ es el conjunto de estados finales. 
\end{itemize}

Se denota $P(A)$, al conjunto de participantes que integran el autómata A, y se define como $P(A) = \left(p \in \mathcal{P} \ | \left(\exists\ \langle p1,p2,c\rangle \in \Sigma_\mathit{Ex}\right)\left(p1=p \lor p2=p\right) \right) $ 
\end{definition}

$\Sigma$ es el conjunto de todas las etiquetas del autómata y lo dividimos en tres subconjuntos disjuntos de la forma $\Sigma_{\mathit{Int}}$, $\Sigma_{\mathit{Ex}}$ y $\Sigma_{\mathit{Buff}}$ que corresponden a las acciones de procesamiento interno, las de comunicación externa con otros participantes y las de comunicación interna vía buffers.

$\Sigma_{Ex}$ es el conjunto de etiquetas correspondientes a la comunicación externa del autómata. Como la comunicación es dirigida punto a punto, cada etiqueta incluye a dos participantes $p_1, P_2 \in \mathcal{P}$ (emisor y receptor) y un canal $c \in \mathcal{C}$. Necesariamente en todas las etiquetas vale que $p_1$ y $p_2$ son distintos y .

$\Sigma_{Buff}$ es el conjunto de etiquetas de acciones de comunicación interna vía buffers. Se especifican las acciones de encolar y desencolar como $b \ll m$, $b \gg m$, respectivamente, donde $b \in B$ y $m$ es el mensaje que se inserta/extrae del mismo.

$\delta$ es el conjunto de transciciones que se compone de $\delta_{Int}$, las acciones internas, $\delta_{Ex}$ las transiciones de comunicación con otros agentes y $\delta_{Buf}$ las acciones de buffer.

$\delta_{Ex}$ se compone de dos tipos de acciones que representan el intercambio de mensajes $m \in \mathcal{M}$ entre dos autómatas a través de un canal de comunicación $c \in \mathcal{C}$. Las acciones de entrada o input representan la recepción de un mensaje de algún proceso externo al correspondiente al autómata, denotadas \textit{In(c,m)}. Las acciones de output o salida que representan el envío de un mensaje y se denotan \textit{Out(c,m)}. Estas acciones son la parte externa de la comunicaión asincrónica. Como los canales son unidireccionales decimos que para todo par de participantes existen una cantidad finita de pares de canales.

$\delta_{Buff}$ es la relación de transición de comunicación interna mediante buffers.\\

Formalmente, F $\subseteq$ Q. Es el conjunto de estados finales, los posibles estados a los que una ejecución llega y es aceptada como válida. Para los autómatas asincrónicos pedimos también que al llegar a este estado final los buffers se encuentren vacíos. Para esto definimos lo siguiente:

\begin{definition}[Configuración instantánea]
Dados $\mathcal{P}$ un conjunto de participantes y $\mathcal{M}$ un conjunto de mensajes. A un autómata con comunicación asincrónica $A_\mathcal{P} = \langle Q, B, \mathcal{C}, \Sigma, \delta, q_0, F\rangle$, definimos:

Una configuración instantánea del autómata como $\langle q, \Omega \rangle \in Q \times \{ [m_b]_{b \in B} | m_b \in \mathcal{M}^* \}$ donde $\mathcal{M}^* = \{ [s_o, \ldots, s_n] \  | \ \forall i \in [0, \ldots, n], s_i \in \mathcal{M} \}$ y $[m_b]_{b \in B} $

\begin{itemize}
\item  q es es el estado actual
\item $\Omega$ es el conjunto de buffers (con sus respectivos contenidos hasta el momento). Decimos que $\Omega$ es de la forma ${ \omega_{b_1}, \omega_{b_2},..., \omega_{b_n} }$ con $b_i \in B$. 
\item Decimos que una configuración es inicial si $q = q_0$ y $ \Omega = \{ [], \ldots, [] \}$
\item Decimos que una configuración es final si $q \in F$ y $ \Omega = \{ [], \ldots, [] \}$
\end{itemize}
\end{definition}

\begin{definition}[Relación de transición entre configuraciones]
\label{def:reltransAFCA}
Sean $A_\mathcal{P} = \langle Q, B, \mathcal{C}, \Sigma, \delta, q_0, F\rangle$, $q_1, q_2 \in Q$, $ m \in \mathcal{M}$ definimos $\vdash$, relación de transición entre configuraciones, como:  
\begin{enumerate}
\item $\langle q_1, \Omega \rangle \vdash \langle q_2, \Omega \rangle \iff \langle q_1,r, q_2 \rangle \in \delta \land r \in \Sigma \setminus \Sigma_\mathit{Buff}$

\item $\langle q_1 ,\{\omega_{b_1},\ldots,\omega_{b_i},\ldots,\omega_{b_n}\} \rangle \vdash \langle q_2,\{\omega_{b_1},\ldots,m : \omega_{b_i},\ldots,b_n\}\rangle \iff \langle q_1, b_i \ll m, q_2 \rangle  \in \delta $ 

\item $\langle q_1, \{\omega_{b_1},\ldots,\omega_{b_i} : m,\ldots,\omega_{b_n}\} \rangle \vdash \langle q_2, \{\omega_{b_1},\ldots,\omega_{b_i},\ldots,\omega_{b_n} \}\rangle \iff \langle q_1, b_i \gg m, q_2 \rangle \in \delta$
\end{enumerate}

\end{definition}

% \begin{definition}[Estados y configuraciones alcanzables]
% \label{def:estadosyconAFCA}
%   Una configuración $c' = \langle q', \Omega' \rangle$ es \emph{alcanzable} desde otra  configuración $c = \langle q, \Omega \rangle$ a través de la \emph{ejecución de la transición} $l$ (escrito $c \overset{l}{\rightarrow} c'$) si existe un $m \in \mathcal{M}$ tal que ocurre una de las siguientes alternativas:
% 	\begin{enumerate}
% 		\item $l = b_i \ll m$, $\langle q, l,  q'\rangle \in \delta$, $\Omega = [\omega_{b_1},\ldots, \omega_{b_i},\ldots,b_n]$ y $\Omega' = [\omega_{b_1},\ldots,m : \omega_{b_i},\ldots,b_n]$
% 		\item $l = b_i \gg m$, $\langle q, l,  q'\rangle \in \delta$, $\Omega = [\omega_{b_1},\ldots,m : \omega_{b_i},\ldots,b_n]$ y $\Omega' = [\omega_{b_1},\ldots, \omega_{b_i},\ldots,b_n]$
% 	\end{enumerate}
% La clausura reflexoEscribimos $ s_1 \overset{t1...tm}{\rightarrow} s_{m+1}$ cuando para algún $s_2,...,s_m, s_1\overset{t}{\rightarrow} s_2...s_m\overset{t_m}{\rightarrow} s_{m+1} $ . El conjunto de configuraciones alcanzables de S es $RS(S) = \{s | s_0 \rightarrow^*S \}$
% \end{definition}

\begin{definition}[Traza de un AFCA] Llamamos traza a una secuencia posible de acciones de un autómata. Se define como una secuencia finita de etiquetas de estado y transición altenradas, que comienza y termina con un estado. Dado un autómata $A_\mathcal{P} = \langle Q, B, \mathcal{C}, \Sigma, \delta, q_0, F\rangle$, una traza tiene las siguientes caracterísitcas

\begin{itemize}
\item Tiene la forma $[q_0, \sigma_1, q_1,...,q_{n-1}, \sigma_n, q_n] $ donde 
\item $q_0$ es el estado inicial del autómata
\item $q_i \in Q$,
\item $ \sigma_i \in \Sigma$ y
\item $\langle q_{i-1}, \sigma_i, q_i \rangle \in \delta $ 
\end{itemize}

El comportamiento de un AFCA es el conjunto de todas las trazas posibles.
\end{definition}




% \begin{definition}[Condición de aceptación] Sea $\sigma$ una secuencia de transiciones sobre $A_{\mathcal{P}}$, decimos que es una ejecución válida si se cumple lo siguiente:

% \centering
% $\sigma \in E(A_{\mathcal{P}}) \iff \exists q_f \in F \ | \  (q_0, \sigma, [[],..,[]]) \vdash^*(q_f, \emptyset, [ [],..,[] ])$ 
% \end{definition}


% Entonces decimos que una ejecución aceptable es aquella que cumple que en una cantidad finita de pasos llegó a un estado final con los buffers vacíos, asegurándonos la ausencia de mensajes huérfanos y deadlocks.


\begin{definition}[Ejecución de un AFCA]. Una ejecución es un $ \tau = \tau_0, \ldots,\tau_n $ donde $\tau_0$ es la configuración inicial, $\langle \tau_i, \tau_{i+1} \rangle \in \vdash$, $\tau_n$ es una configuración final y $ \forall i \in [1, \ldots, n], \tau_i$ es una configuración válida. Una ejecución es una configuración que no entra en deadlock, mensajes huérfanos ni receptor no especificado.
\end{definition}

\begin{ejemplo}[Ejemplo de un Autómata Finito de Comunicación Asincrónica]
\label{ex:AFCA}
Considere el AFCA $S= \langle \{q_0,..,q_5\},[b_0],\{sr_1,sr_2,rs_1\}, \Sigma_A, \delta_A, q_0, \{q_4\} \rangle$. Donde $\Sigma =\{wait, b_0 \ll m_1,b_0 \gg m_1, Out(sr_1,a),Out(sr_2,c), In(rs_1,b),In(rs_2,d) \}$
\begin{center}
\begin{tikzpicture}[->, thick]
 \node[state,initial] (q_0)   {$q_0$}; 
 \node[state] (q_1) [right= 1.5cm of q_0 ] {$q_1$};
 \node[state] (q_5) [right= 1.8cm of q_1 ] {$q_5$};
 \node[state] (q_2) [below= of q_0 ] {$q_2$};
 \node[state] (q_3) [right= 1.5cm of q_2 ] {$q_3$};
 \node[state,accepting] (q_4) [right= 1.8cm of q_3 ] {$q_4$};
 \draw[]        
        (q_0) edge[above] node{$b_0 \ll m_1$} (q_1)
        (q_0) edge[left] node{$Out(sr_1,a)$} (q_2)
        (q_1) edge[above] node{$Out(sr_2,c)$} (q_5)
        (q_2) edge[above] node{$In(rs_1,b)$} (q_3)
        (q_2) edge[bend right, below] node{$wait$} (q_4)
        (q_3) edge[above] node{$In(rs_2,d)$} (q_4)
        (q_5) edge[right] node{$b_0 \gg m_1$} (q_4)
        ;
\end{tikzpicture}
\end{center}
En el ejemplo podemos ver los tres tipos de transiciones que los AFCA pueden realizar. 

\begin{itemize}
    \item Transiciones de buffer interno: $b_0 \ll m_1$ y $b_0 \gg m_1$
    \item Transiciones de comunicación externa: $Out(sr_1,a),Out(sr_2,c), In(rs_1,b)$ y $In(rs_2,d) $
    \item Transiciones internas: $wait$
\end{itemize}
\end{ejemplo}

Como la idea es que los autómatas con comunicación asincrónica representen procesos, o servicios, que pueden formar parte de un sistema más grande, necesitamos definir la operación de composición. 

\begin{definition}[Compatibilidad]
Dados $\mathcal{P}$ un conjunto de participantes, $\mathcal{M}$ un conjunto de mensajes y $E_\mathcal{P} = \{E_{p} \ | p \in \mathcal{P}\}$ un conjunto finito de autómatas finitos de comunicación asincrónica $E_{p_i}= \langle Q_{p_i}, B_{p_i}, \mathcal{C}_{p_i}. \Sigma_{p_i}, \delta_{p_i}, q_{0{p_i}}, F_{p_i} \rangle$, diremos que $E_\mathcal{P}$ es \emph{compatible} si se satisface que para todos $p_i, p_j \in \mathcal{P}$, $\Sigma_{p_i} \cap \Sigma_{p_j}= \emptyset$ y $B_{p_i} \cap B_{p_j} = \emptyset$.
\end{definition}

\begin{definition}[Composición]
\label{def:composicion}
Dados $\mathcal{P}$ un conjunto de participantes, $\mathcal{M}$ un conjunto de mensajes y $E_\mathcal{P} = \{E_{p} \ | p \in \mathcal{P}\}$ un conjunto finito de autómatas finitos de comunicación asincrónica $E_{p_i}= \langle Q_{p_i}, B_{p_i}, \mathcal{C}_{p_i}. \Sigma_{p_i}, \delta_{p_i}, q_{0{p_i}}, F_{p_i} \rangle$ compatible. Luego, definimos la composición $E_\mathcal{P} = ||_{1..n} E_{p_i}$ como sigue:
\begin{itemize}
    \item $Q_\mathcal{P}= \Pi_{p_i \in \mathcal{P}} Q_{p_i}$ (i.e. el conjunto de estados de la composición es el producto cartesiano de los estados de los autómatas componentes),    
    \item $B_\mathcal{P} = \bigcup_{p_i \in \mathcal{P}} B_{p_i} \cup \{ \omega_c \ | \ \mbox{there exists } p_i, p_j \in \mathcal{P} \mbox{ such that } c \in \mathcal{C}_{p_i} \cap \mathcal{C}_{p_j} \}$ (i.e. el conjunto de nombres de buffers del autómata resultante se compone de los buffers de cada autómata participante en la composición, junto con uno nuevo por cada canal compartido entre cada par de autómatas, dichos canales son aquellos mediante los cuales los autómatas a componer intercambian mensajes),    
    \item $\mathcal{C}_\mathcal{P} = \bigcup_{p_i \in \mathcal{P}} \mathcal{C}_{p_i} \setminus \{ c_k \ | \mbox{for all } p_i, p_j \in \mathcal{P}, c \in \mathcal{C}_{p_i} \cap \mathcal{C}_{p_j} \}$,    
    \item $\Sigma_\mathcal{P} = \Sigma_{\mathcal{P}\mathit{Int}} \cup  \Sigma_{\mathcal{P}\mathit{Ex}} \cup \Sigma_{\mathcal{P}\mathit{Buff}}$ tal que:    
    \begin{inparaenum}[1)]
        \item $\Sigma_{\mathcal{P}\mathit{Int}} = \bigcup_{p_i \in \mathcal{P}} \Sigma_{p_i\mathit{Int}}$,
        \item $\Sigma_{\mathcal{P}\mathit{Ex}} = \bigcup_{p_i \in \mathcal{P}} \Sigma_{p_i\mathit{Ex}} \setminus \Sigma_{p_i \mapsto p_j}$ y 
        \item $\Sigma_{\mathcal{P}\mathit{Buff}} = \bigcup_{p_i \in \mathcal{P}} \Sigma_{p_i\mathit{Buff}} \cup \Sigma_{p_i \mapsto p_j}$,
    \end{inparaenum}
donde $\Sigma_\mathit{p_i \mapsto p_j} =\{ \langle p_i,p_j,c \rangle \ | \ p_i, p_j \in \mathcal{P}, \  c \in \mathcal{C}_{p_i} \cap \mathcal{C}_{p_j}\}$
\item $\delta_{\mathcal{P}} = \delta_{\mathcal{P}\mathit{Int}} \cup \delta_{\mathcal{P}\mathit{Ex}} \cup \delta_{\mathcal{P}\mathit{Buff}}$ tal que:
\begin{enumerate}
\item $\delta_{\mathcal{P}\mathit{Int}} \subseteq Q_{\mathcal{P}} \times \Sigma_{\mathcal{P} Int} \times Q_{\mathcal{P}}$  es la relación de transición interna, tal que que se satisface la siguiente fórmula:
$$(\forall \langle q, \sigma, q' \rangle \in \delta_{\mathcal{P}\mathit{Int}})(\exists p_i \in \mathcal{P})(\exists q_{p_i}, q'_{p_i} \in Q_{p_i}, \sigma \in \Sigma_{p_i})(\langle q_{p_i}, \sigma, q'_{p_i} \rangle \in \delta_{p_i Int})$$

\item $\delta_{\mathcal{P}\mathit{Ex}} \subseteq Q_{\mathcal{P}} \times \Sigma_{\mathcal{P}\mathit{Ex}} \times Q_{\mathcal{P}}$ es la relación de transición externa, tal que que se satisface la siguiente fórmula: $$(\forall \langle q, \sigma, q' \rangle \in \delta_{\mathcal{P}\mathit{Ex}}) (\exists p_i \in \mathcal{P})(\exists q_{p_i},q'_{p_i} \in Q_{p_i}, \sigma \in \Sigma_{p_i}) (\langle q_{p_i},\sigma,q'_{p_i} \rangle \in \delta_{p_i Ex})$$

\item $\delta_{\mathcal{P}\mathit{Buff}} \subseteq Q_{\mathcal{P}} \times \Sigma_{\mathcal{P} \mathit{Buff}} \times Q_{\mathcal{P}}$ es la relación de transición por comunicación asincrónica interna tal que se satisface la siguiente fórmula:
$$
\begin{array}{l}
\langle q, \sigma, q' \rangle \in \delta_{\mathcal{P}\mathit{Buff}} \mbox{ sii } (q = \langle q_k \rangle_{p_k \in \mathcal{P}} \land q' = \langle q'_{k'} \rangle_{p_{k'} \in \mathcal{P}} \land\\
\qquad
\begin{array}{l}
((\exists p_i \in \mathcal{P})((\forall p_k \in \mathcal{P})(p_k \neq p_i \implies q_k = q'_k) \land \\
\qquad(\exists \omega_b \in B_{p_i}, \omega_b \{\gg, \ll\} m \in \Sigma_{{p_i} \mathit{Buff}})(\langle q_i, \omega_b \{\gg, \ll\} m, q'_i \rangle \in \delta_{p_i \mathit{Buff}}) \land\\
\qquad\qquad \sigma = \omega_b \{\gg, \ll\} m)\lor\\
(\exists p_i, p_j \in \mathcal{P})(p_i \neq p_j \land (\forall p_k \in \mathcal{P})((p_k \neq p_i \land p_k \neq p_j) \implies q_k = q'_k) \land \\ 
\qquad ((\exists \mathit{Out}(c_{{p_i, p_j}_n}, m) \in \Sigma_{{p_i} \mathit{Ex}})(\langle q_i, \mathit{Out}(c_{{p_i, p_j}_n}, m), q'_i \rangle \in \delta_{p_i  \mathit{Ex}} \land\\
\qquad\qquad \sigma = \omega_{c_{{p_i p_j}_n}} \ll m) \lor \\ 
\qquad (\exists \mathit{In}(c_{p_j,{p_i}_n}, m) \in \Sigma_{{p_i} \mathit{Ex}})(\langle q_i, \mathit{In}(c_{p_j,{p_i}_n}, m), q'_i \rangle \in \delta_{p_i \mathit{Ex}} \land\\
\qquad\qquad \sigma = \omega_{c_{p_j,{p_i}_n}} \gg m))))
\end{array}
\end{array}
 $$
 
Además vale que $\forall q_i,q_j \in Q_\mathit{\mathcal{P}}, \ m \in \mathcal{M} \ |	 \ \langle q_i, [\omega_\mathit{b}]_{b \in B_\mathcal{P}} \gg m, q_j \rangle$

$\langle q_i, [\omega_\mathit{b}]_{b \in B_\mathcal{P}} \gg m, q_j \rangle \in \delta_\mathit{\mathcal{P}Buff} \iff \exists q'_i, q'_j \in Q_\mathit{\mathcal{P}}$ tal que $\ \langle q'_i, [\omega_\mathit{b}]_{b \in B_\mathcal{P}} \ll m, q'_j \rangle \in \delta_\mathit{\mathcal{P}Buff}$ y se cumple una de las siguientes condiciones:
\begin{enumerate}
\item $q'_j = q_i$ 
\item $\exists \sigma \in \Sigma_\mathit{\mathcal{P}}$ tal que $\langle q'_j, \sigma, q_i \rangle \in \delta_\mathit{\mathcal{P}} \land \sigma \neq [\omega_\mathit{b}]_{b \in B_\mathcal{P}} \gg m $
\item $\exists s=[q'_j,...,q_i] $ donde s es una secuencia finita de estados tales que $s[0]= q'_j, s[n-1]=q_i$ y sea $ 0 \ll x \ll$ n-1,$s[x] \in Q_\mathit{\mathcal{P}}$ y $\forall s[x], s[x+1], \exists \sigma \in \Sigma_\mathit{\mathcal{P}} \land \langle s[x], \sigma, s[x+1] \rangle \in \delta_\mathit{\mathcal{P}} $  
\end{enumerate} 
\end{enumerate}
\item $q_0 = \langle {q_0}_{p_1}, \ldots, {q_0}_{p_n} \rangle$, y
\item $F_{\mathcal{P}} = \bigcup_{1..n} F_{p_i}$.
\end{itemize}
\end{definition}


Cada autómata es un sistema independiente que cumple una función (o una serie de funciones), y se relaciona con otros a través del envío de mensajes. Para asegurarnos esto pedimos que para todo par de autómatas a componerse tengan conjuntos de etiquetas de acciones que sean disjuntos. 

Del mismo modo, cada autómata tiene su propio conjunto de buffers que pedimos sean disjuntos, para distinguir la comunicación interna de cada autómata componente de la que ocurra entre componentes o con participantes externos. Para modelar la comunicación interna entre componentes, agregamos dos buffers por cada par de integrantes, uno para cada sentido de la comunicación, de $E_{p_i}$ a $E_{p_j}$ y viceversa.

Como los conjuntos de acciones son disjuntos podemos decir que las acciones internas, de comunicación externa y de buffer, de cada componente se preservan, siempre y cuando tenga sentido con la composición de estados. Hay un caso particular que se da cuando ocurren envíos de mensaje de un autómata componente a otro. En ese caso dado que ambos ahora son parte un mismo autómata, la comunicación pasa a ser envío de mensajes interno. Para representar este tipo de comunicación es que utilizamos buffers. De este modo el intercambio que antes era $\mathit{In} (c_{p_i, p_{j_n}},m)$ y $\mathit{Out}(c_{p_j, p_{i_n}},m)$ ahora es $\omega_{c_{p_i p_{j_n}}}  \ll m$ y $b_{c_{p_j, p_{i_n}}}  \gg  m$, donde $ \omega_{c_{p_i, p_{j_n}}}, \omega_{c_{p_j, p_{i_n}}} \in \mathcal{C}_{\mathcal{P}} $ son los buffers exclusivos del autómata compuesto. 

Al componer autómatas la comunicación que antes era externa y ahora es de buffer puede generar problemas. Puntualmente pueden aparecer transiciones de consumo de un buffer (que antes eran envío de mensajes) donde antes no había. De este modo pueden aparecer secuencias de estados y transiciones donde se consume un mensaje antes de que este sea depositado en el buffer correspondiente. Para esto pedimos que $\delta_\mathit{\mathcal{P}Buff}$ cumpla con una condición especial. Solo pueden haber transiciones de consumo saliendo de un estado si en alguna secuencia de acciones que termina en ese estado, hay transiciones de producción (es decir se encola un mensaje en el buffer).

\begin{ejemplo}[Composición de Autómatas Finitos de Comunicación Asincrónica]
\label{ex:Composicion}
%%%%%%%%%%
% Ejemplo viejo
%%%%%%%%%%
%
%Considere los siguientes AFCA:
%\begin{itemize}
%    \item $S= \langle \{q_0,..,q_4\},[b_0],\{sr_1,rs_1\}, \Sigma_S, \delta_S, q_0, \{q_4\} \rangle$, con $\Sigma_S =\{wait, b_0 \ll m_1,b_0 \gg m_1, In(rs_1,b), Out(sr_1,c) \}$
%\begin{center}
%\begin{tikzpicture}[->, thick][scale=0.33]
% \node[state,initial] (q_0)   {$q_0$}; 
% \node[state] (q_1) [right= 1.5cm of q_0 ] {$q_1$};
% \node[state] (q_2) [right = 1.5 of q_1 ] {$q_2$};
% \node[state] (q_3) [below=  of q_1 ] {$q_3$};
% \node[state,accepting] (q_4) [left= 2cm of q_3 ] {$q_4$};
% \draw[]        
%        (q_0) edge[above] node{$In(rs_1,b)$} (q_1)
%        (q_1) edge[above] node{$b_0 \ll m_1$} (q_2)
%        (q_1) edge[left] node{$wait$} (q_3)
%        (q_2) edge[right] node{$\ b_0 \gg m_1$} (q_3)
%        (q_3) edge[above] node{$Out(sr_1,c)$} (q_4)
%        ;
%\end{tikzpicture}
%\end{center}
%        \item $R=\langle \{q_0,q_1,q_2\},[],\{sr_1,rs_1\}, \Sigma_R, \delta_R, q_0, \{q_4\} \rangle$, con $\Sigma_R = \{Out(rs_1,b), In(sr_1,c)\}$, y su composición
%\begin{center}
%\begin{tikzpicture}[->, thick][scale=0.33]
% \node[state,initial] (q_0)   {$q_0$}; 
% \node[state] (q_1) [below= of q_0 ] {$q_1$};
%  \node[state,accepting] (q_2) [below= of q_1] {$q_2$};
% \draw[]        
%        (q_0) edge[right] node{$Out(rs_1,b)$} (q_1)
%        (q_1) edge[right] node{$In(sr_1,c)$} (q_2)
%        ;
%\end{tikzpicture} 
%\end{center}
%    \item $\RS =\langle Q_{RS}, [b_0, sr_1, rs_1], \Sigma_{RS}, \delta_{RS}, q_{00}, \{q_{42}\} \rangle$, con $\Sigma_{RS} =\{wait, b_0 \ll m_1,b_0 \gg m_1, rs_1 \ll b, rs_1 \gg b, sr_1 \ll c, sr_1 \gg c) \}$
%\begin{center}
%\begin{tikzpicture}[->, thick][scale=0.33]
% \node[state,initial] (q_00)   {$q_{00}$}; 
%  \node[state] (q_01) [right= 1.8cm of q_00 ] {$q_{01}$};
%  \node[state] (q_11) [right= 1.8cm of q_01 ] {$q_{11}$};
%  \node[state] (q_21) [right= 1.8cm of q_11 ] {$q_{21}$};
%  \node[state] (q_31) [below= of q_21 ] {$q_{31}$};
%  \node[state] (q_41) [left= 1.8cm of q_31 ] {$q_{41}$};
%  \node[state,accepting] (q_42) [left= 1.8cm of q_41] {$q_{42}$};
%  \draw[]        
%         (q_00) edge[above] node{$rs_1 \ll b$} (q_01)
%         (q_01) edge[above] node{$rs_1 \gg b$} (q_11)
%         (q_11) edge[above] node{$b_0 \ll m$} (q_21)
%         (q_11) edge[right] node{$wait$} (q_31)
%         (q_21) edge[right] node{$b_0 \gg m$} (q_31)
%         (q_31) edge[above] node{$sr \ll c$} (q_41)
%         (q_41) edge[above] node{$sr \gg c$} (q_42)
%       ;
%\end{tikzpicture}
%\end{center}
%\end{itemize} 
%
%%%%%%%%%%

\cnote{Explicar cómo los autómatas ejemplifican la comunicación referenciando la Fig.~\ref{a_i}.}

\begin{figure}[H]
\begin{center}
\begin{tikzpicture}[->, thick]
 \node[state,initial] (q_0)   {$q_0$}; 
 \node[state] (q_1) [right= 1.5cm of q_0 ] {$q_1$};
 \node[state] (q_2) [below= of q_0 ] {$q_2$};
 \node[state] (q_3) [right= 1.5cm of q_2 ] {$q_3$};
 \node[state, accepting] (q_4) [below= of q_3 ] {$q_4$};
	
 \draw[]        
        (q_0) edge[above] node{out(pr,a)} (q_1)
        (q_0) edge[left] node{in(sp,b)} (q_2)
        (q_1) edge[right] node{in(sp,b)} (q_3)
        (q_2) edge[above] node{out(pr,a)} (q_3)
        (q_3) edge[right] node{$int_p$} (q_4)
        ; 
\end{tikzpicture} 
\ 
\begin{tikzpicture}[->, thick]
 \node[state,initial] (q_0)   {$q_0$}; 
 \node[state] (q_1) [below= of q_0 ] {$q_1$};
 \node[state, accepting] (q_2) [below= of q_1 ] {$q_2$};
 \draw[]        
        
        (q_0) edge[right] node{$int_r$} (q_1)
        (q_1) edge[right] node{in(pr,a)} (q_2)
        ;
\end{tikzpicture} 
\ 
\begin{tikzpicture}[->, thick]
 \node[state,initial] (q_0)   {$q_0$}; 
 \node[state] (q_1) [below= of q_0 ] {$q_1$};

 \draw[]        
        
        (q_0) edge[right] node{out(sp,b)} (q_1)
        
        ;
\end{tikzpicture} 
\end{center}
\caption{Autómatas finitos de comunicación asincrónica $\mathcal{A}_1$, $\mathcal{A}_2$, $\mathcal{A}_3$.}
\label{a_i}
\end{figure}

\begin{figure}[H]
\begin{center}
\begin{tikzpicture}[->, thick]
 \node[state,initial] (q_0)   {$q_{000}$}; 
 \node[state] (q_1) [right= 1.5cm of q_0 ] {$q_{100}$};
 \node[state] (q_2) [right = 1.5 of q_1 ] {$q_{110}$};
 \node[state] (q_3) [right = 1.5 of q_2 ] {$q_{120}$};
  \node[state] (q_4) [right = 1.5 of q_3 ] {$q_{121}$};
  \node[state] (q_5) [below=  of q_0] {$q_{001}$};
  \node[state] (q_6) [right = 1.5cm of q_5] {$q_{201}$};
  \node[state] (q_7) [right = 1.5 of q_6 ] {$q_{301}$};
  \node[state] (q_8) [right = 1.5 of q_7 ] {$q_{311}$};
  \node[state] (q_9) [right = 1.5 of q_8 ] {$q_{321}$};
 \node[state,accepting] (q_10) [right= 1.5cm of q_9 ] {$q_{421}$};
 \draw[]        
        (q_0) edge[above] node{$PR \ll a$} (q_1)
        (q_1) edge[above] node{$int_r$} (q_2)
		(q_2) edge[above] node{$PR \gg a$} (q_3)
        (q_3) edge[above] node{$SP \ll b$} (q_4) 
        (q_4) edge[right] node{$SP \gg b$} (q_9)
        (q_0) edge[left] node{$SP \ll b$} (q_5)
        (q_5) edge[below] node{$SP \gg b$} (q_6)
        (q_6) edge[below] node{$PR \ll a$} (q_7)
        (q_7) edge[below] node{$int_r$} (q_8)
        (q_8) edge[below] node{$PR \gg a$} (q_9)
        (q_9) edge[below] node{$int_p$} (q_10)
        ;
\end{tikzpicture}
\end{center}
\caption{Composición de los AFCAs de la Fig.~\ref{a_i}.}
\label{A}
\end{figure}



\end{ejemplo}

\subsection*{AFCAs y composición parcial}
\cnote{Escribir cómo los AFCAs posibilitan composición parcial y explicar referenciando la Fig.~\ref{A-parcial}.}

\cnote{Componer $\mathcal{A}_1$ y $\mathcal{A}_2$ y dejar $\mathcal{A}_3$ como está.}

\begin{figure}[H]
\begin{center}
\begin{tikzpicture}[->, thick]
 \node[state,initial] (q_0)   {$q_{000}$}; 
 \node[state] (q_1) [right= 1.5cm of q_0 ] {$q_{100}$};
 \node[state] (q_2) [right = 1.5 of q_1 ] {$q_{110}$};
 \node[state] (q_3) [right = 1.5 of q_2 ] {$q_{120}$};
  \node[state] (q_4) [right = 1.5 of q_3 ] {$q_{121}$};
  \node[state] (q_5) [below=  of q_0] {$q_{001}$};
  \node[state] (q_6) [right = 1.5cm of q_5] {$q_{201}$};
  \node[state] (q_7) [right = 1.5 of q_6 ] {$q_{301}$};
  \node[state] (q_8) [right = 1.5 of q_7 ] {$q_{311}$};
  \node[state] (q_9) [right = 1.5 of q_8 ] {$q_{321}$};
 \node[state,accepting] (q_10) [right= 1.5cm of q_9 ] {$q_{421}$};
 \draw[]        
        (q_0) edge[above] node{$PR \ll a$} (q_1)
        (q_1) edge[above] node{$int_r$} (q_2)
		(q_2) edge[above] node{$PR \gg a$} (q_3)
        (q_3) edge[above] node{$SP \ll b$} (q_4) 
        (q_4) edge[right] node{$SP \gg b$} (q_9)
        (q_0) edge[left] node{$SP \ll b$} (q_5)
        (q_5) edge[below] node{$SP \gg b$} (q_6)
        (q_6) edge[below] node{$PR \ll a$} (q_7)
        (q_7) edge[below] node{$int_r$} (q_8)
        (q_8) edge[below] node{$PR \gg a$} (q_9)
        (q_9) edge[below] node{$int_p$} (q_10)
        ;
\end{tikzpicture}
\end{center}
\caption{Composición parcial de los AFCAs de la Fig.~\ref{a_i}.}
\label{A-parcial}
\end{figure}



\section{Aspectos comunicacionales de los AFCA}
Una de las desventajas principales de utilizar CFSM como lenguaje de especificación de en Service Oriented Computing es que la naturaleza del binding en SOC choca con el método de establecer si un conjunto de participantes pueden interacturar libres de errores de comunicación. En SOC el descubrimiento y binding de servicios se realiza por demanda, eso significa que a pesar de que se necesiten muchos servicios para resolver una tarea, se obtienen uno por uno según se van necesitando. Por el otro lado las CFSM requieren el conjunto total de participantes de un protocolo para poder determinar si el protocolo puede llevarse a cabo sin errores. La idea de los AFCA es intentar reducir la distancia entre ambas naturalezas. Por lo tanto una de las tareas principales que queremos lograr con AFCAs es obtener, vía una proyección, la inferfaz de comunicación en la forma de una mCFSM.
Al principio de esta sección introdujimos los AFCA para modelar el comportamiento deseado en SOC de binding incremental. Ahora nos queremos enfocar en los aspectos puramente comunicacionales, es decir la interacción entre distintos participantes. Para hacer énfasis en estos aspectos existen las CFSMs, pero los AFCA introducen con la composición la posibilidad de que haya múltiples canales de comunicación entre dos participantes. En esta sección abordamos dicha problemática introduciendo las Multichannel Communicating Finite State Machines \cite[Def.~82]{vissani:phdthesis}. Se definen, del mismo modo que una CFSM sobre $\mathcal{M}$ de la siguiente manera.

\begin{definition}[Multichannel Communicating Finite State Machine - mCFSM]\label{mCFSM} Una multichannel communicating finite state machine (mCFSM) sobre $\mathcal{M}$ es una un sistema finito de transición $(Q, \mathcal{C}, q_0, \mathcal{M}, \delta)$ donde
\begin{itemize}
  \item $Q$ es un conjunto finito de estados
  \item $\mathcal{C} = \{ pq_n \mid pq \in \mathcal{P}^2, n \in \mathbb{N}, p \not= q\}$ es un conjunto de canales
  \item $q_0 \in Q$ es el estado inicial;
  \item $\delta \subseteq Q \times (\mathcal{C} \times \{!,?\} \times \mathcal{M}) \times
    Q$ es un conjunto finito de \emph{transiciones}.
  \end{itemize}

Un CS es un mapa S que asigna un mCFSM S(p) a cada $p \in \mathcal{P}$. Escribimos $q \in S(p)$ cuando q es un estado de la máquina S(p) y $\tau$ es una transición de S(p).\end{definition}

Al igual que antes la semántica de un CS se obtiene considerando configuraciones. Estas configuraciones son iguales a las de las CFSM puras con la salvedad que ahora los canales no están restringidos a un único par entre cada par de participantes.

 
% * <ivissani@gmail.com> 2018-01-28T13:14:43.225Z:

% De esta manera lo que tenés es una enumeración de canales entre cada par de participantes. Lo único que hay que hacer es ajustar la semántica (las reglas) a esta nueva notación para que sea coherente.
% 
% Notar que la cantidad de canales existentes entre cada par de participantes ahora es infinita (los naturales) PERO la cantidad de ellos que vas a usar siempre va a ser finita porque de acuerdo a la notación que tenemos hay que fijar qué canal se usa en cada envío/recepción de mensaje y, por lo tanto, el hecho de que EXISTAN infinitos canales no afecta a la semántica ni al poder expresivo de las CFSMs
% 

Multichannel CFSMs  con un único canal para cada par ordenado de participantes son equivalentes a CFSMs puras.

\begin{definition}[Interfaz de comunicación de un AFCA]\label{def:interfazAFCA} Es la mCFSM que resulta de aplicarle el siguiente procedimiento:
\begin{enumerate}
    \item Transformamos toda acción interna y de buffer en una transición $\epsilon$. Esto resulta en un autómata no determinístico con transiciones $\epsilon$ al que llamamos $A'$
    \item $\forall q \in Q_{A'}$ llamamos $[q]$ a la clase de equivalencia tal que $[q]=\{q' \in Q_{A'}\ |\ q \xrightarrow{\epsilon^*} q'\}$. Decimos que dos clases $[q]$ y $[q']$ están relacionadas $[q] \xrightarrow{\sigma} [q'']$, con $\sigma \in \Sigma_{Ex}$, sii $\exists q^1 \in [q], q^2 \in [q'']$ tales que $q^1 \xrightarrow{\sigma} q^2 \in \delta$
    \item $[Q]= \{[q] \ | \ q, q' \in Q \land [q] \xrightarrow{\sigma} [q'] \lor [q'] \xrightarrow{\sigma} [q] \}$ es el conjunto de clases de equivalencia de $Q$ que están relacionadas. 
    \item $[q_0]=q_0$
    \item $[\delta]= \{[q] \xrightarrow{\sigma} [q']\ |\ [q], [q'] \in [Q] \land \sigma \in \Sigma_{EX}\}$, dos estados $[q],[q'] \in [Q]$ están relacionados en $[\delta]$ si son clases de equivalencia relaicionadas. 
\end{enumerate}
\end{definition}

\begin{ejemplo}[Interfaz de comunicación de un AFCA] 
\label{def:interfazcom}
%%%%%%%%%%
% Ejemplo viejo
%%%%%%%%%%
%
%Consiere el siguiente autómata finito de comunicación asincrónica:
%\begin{center}
%        \begin{tikzpicture}[->, thick]
%            \node[state,initial] (q_0)   {$q_0$}; 
%            \node[state] (q_1) [right= 1.8cm of q_0 ] {$q_1$};
%            \node[state, accepting] (q_3) [right= 1.5cm of q_1 ] {$q_3$};
%            \node[state] (q_5) [right= 1.8cm of q_3 ] {$q_5$};
%            \node[state] (q_2) [below= 1.5cm of q_0 ] {$q_2$};
%            \node[state] (q_4) [right= 1.8cm of q_2 ] {$q_4$};
%    
%            \draw[]        
%            (q_0) edge[above] node{$out(pr_1,a)$} (q_1)
%            (q_1) edge[above] node{$in(rp_2,d)$} (q_3)
%            (q_0) edge[left] node{$in(rp_1,b)$} (q_2) 
%            (q_2) edge[below] node{$out(pr_2,c)$} (q_4)
%            (q_4) edge[left] node{$intr$} (q_3)
%            (q_3) edge[bend right,below] node{$b_1 \ll m_1$} (q_5)
%            (q_5) edge[bend right,above] node{$b_1 \gg m_1$} (q_3)
%        ;
%        \end{tikzpicture}
%        \end{center}
%A continuación, presentamos su interfaz de comunicación formalizada como una multichannel communicating finite state machine:
%\begin{center}
%    \begin{tikzpicture}[->, thick]
%            \node[state,initial] (q_0)   {$[q_0]$}; 
%            \node[state] (q_1) [right= 1.5cm of q_0 ] {$[q_1]$};
%            \node[state] (q_2) [below= 1.5cm of q_0 ] {$[q_2]$};
%            \node[state] (q_3) [right= 1.5cm of q_2 ] {$[q_3]$};
%
%            \draw[]        
%            (q_0) edge[above] node{$pr_1!a$} (q_1)
%            (q_0) edge[right] node{$rp_1?b$} (q_2)
%            (q_1) edge[right] node{$rp_2?d$} (q_3)
%            (q_2) edge[above] node{$pr_2!c$} (q_3)
%        ; 
%        \end{tikzpicture}     
%\end{center}
%
%%%%%%%%%%

\cnote{Explicar el ejemplo referenciando la Fig.~\ref{m_i}.}

\begin{figure}[H]
\begin{center}
\begin{tikzpicture}[->, thick]
 \node[state,initial] (q_0)   {$q_0$}; 
 \node[state] (q_1) [right= of q_0 ] {$q_1$};
  \node[state] (q_2) [below= of q_0 ] {$q_2$};
 \node[state] (q_3) [right= of q_2 ] {$q_3$};

 \draw[]        
        (q_0) edge[above] node{pr!a} (q_1)
        (q_0) edge[right] node{sp?b} (q_2)
        (q_1) edge[right] node{sp?b} (q_3)
        (q_2) edge[above] node{pr!a} (q_3)
        ; 
\end{tikzpicture} 
\
\begin{tikzpicture}[->, thick]
 \node[state,initial] (q_0)   {$q_0$}; 
 \node[state] (q_1) [below= of q_0 ] {$q_2$};

 \draw[]        
        (q_0) edge[right] node{pr?a} (q_1)
        ;
\end{tikzpicture} 
\ 
\begin{tikzpicture}[->, thick]
 \node[state,initial] (q_0)   {$q_0$}; 
 \node[state] (q_1) [below= of q_0 ] {$q_1$};
 
 \draw[]        
        (q_0) edge[right] node{sp!b} (q_1)
        ;
\end{tikzpicture} 
\end{center}
\caption{mCFSM $M_1$, $M_2$ y $M_3$, correspondientes a los AFCAs del Subfig.~\ref{a_i}.}
\label{m_i}
\end{figure}

\end{ejemplo}


En la sección 2.3 explicamos como funciona la condición de Generalized Multiparty Compatibility para CFSMs. Necesitamos ver que la misma condición es aplicable a las Multichannel CFSMs, para esto la solución más práctica que se encontró es ver que podemos emular una Multichannel CFSM con una CFSM pura. Mostrando esa equivalencia vemos que si podemos aplicar GMC al emulador, la condición aplica al emulado. Dado que un multichannel CS es libre de errores de comunicación si y solo si el sistema emulado también lo es. El procedimiento consiste en generar un nuevo participante para cada canal entre otros dos participantes. Este nuevo participante es un simple repetidor de mensajes de uno de los participantes a otro. De este modo una comunicación multichannel entre un par de participantes es reemplazada por múltiples comunicaciones de un canal con un repetidor en el medio. El aspecto clave de esta emulación es que preserva el orden de los mensajes.

\begin{definition}[Sistema emulado, \cite{vissani:phdthesis}, Def.~84] Dado un multichannel CS $(M_p)_{p \in \mathcal{P}}$, agrandamos el conjunto $ \mathcal{P}$ agregando un participante adicional por cada canal en el sistema original $\mathcal{P}'=\mathcal{P} \cup \bigcup_{p \in \mathcal{P}} \{p^{pq_n} \ | \ pq_n \in \mathcal{C}_p \} \cup \bigcup_{p \in \mathcal{P}} \{p^{qp_n} \ | \ qp_n \in \mathcal{C}_p \}$. Cada participante nuevo $p^{sr_n} \in \bigcup_{p \in \mathcal{P}} \{p^{pq_n} \ | \ pq_n \in \mathcal{C}_p \} \cup \bigcup_{p \in \mathcal{P}} \{p^{qp_n} \ | \ qp_n \in \mathcal{C}_p \}$ se define con la siguiente mCFSM:
\begin{itemize}
\item $Q_{p^{sr_n}} = \{q_0\} \cup \bigcup_{m \in \mathcal{M}} \{q_m\}$
\item $\mathcal{C}_{p^{sr_n}} = \{sp_n^{sr_n}, p^{sr_n}r_n,  p^{sr_n}s_n, rp_n^{sr_n} \}$
\item $q_{0_{p^{sr_n}}} = q_0 $
\item $ \delta_{p^{sr_n}} = \bigcup_{m \in \mathcal{M}} \{ (q_0,sp_n^{sr_n}?m,q_m),(q_m,p^{sr_n}r_n!m,q_0)  \} $
\end{itemize}
Cada viejo participante $q \in \mathcal{P}$ se reemplaza por un nuevo participante $q'$ donde:
\begin{itemize}
\item $\mathcal{C}_{q'} = \{qp_n^{qr_n} \ | \ qr_n \in \mathcal{C}_{q} \} \cup \{p^{sq_n}q_n \ | \ sq_n \in \mathcal{C}_{q} \}$ 
\item $ \delta_{q'} = \bigcup_{m \in \mathcal{M}} \{ (q, qp_n^{qr_n}!m, q') \ | \ (q, qr_n!m, q') \in \delta_q \} \cup \{q, p^{sq_n}q_n?m,q' \ | \ (q, sq_n?m, q') \in \delta_q \}$

\end{itemize}
\end{definition} 

Queda claro que si bien transformamos cada canal (buffer) en dos canales nuevos, el orden de los mensajes está garantizado como resultado del modo en que $\delta$ está definido para los repetidores. Nótese que los repetidores respetan ese orden porque al consumir un mensaje desde el canal de 'entrada' (el canal que sirve para recibir mensajes del emisor original) la única acción posible del repetidor es mandar el mensaje a través del canal de salida (el canal que sirve para enviar mensajes al receptor original). Como todo canal es FIFO y los repetidores se comportan del mismo modo, el orden entre dentro de cada canal queda garantizado. Por oro lado, no hay garantía respecto al orden de la comunicación entre canales del modelo original. Por lo tanto podemos decir que los repetidores no introducen más concurrencia que la que estaba en el modelo original y debido a eso los dos sistemas, a pesar de no ser bisimilares y no tener las mismas trazas, son equivalentes respecto a la ausencia de deadlock, recepciones no especificadas y mensajes huérfanos. Esta observación es importante porque provee una forma de chequear multichannel communication systems recurriendo nuevamente al chequeo de GMC del sistema emulado.

\begin{prop} El sistema emulado preserva los errores comunicacionales

\begin{proof} Primero notemos que los repetidores no pueden generar errores de comunicación dado que en el estado inicial pueden recibir el rango completo de mensajes y luego de ocnsumir un mensaje su única acción posible es reenviarlo. Recordemos también que el orden de los mensajes se preserva. Esto significa que si en el sistema original los mensajes $ m $ y $m'$ fueron enviados en ese orden sobre el canal $sr$ en el sistema emulado se envían en el mismo orden sobre el canal $sp^{sr}$ y por lo tanto son enviados en ese mismo orden a través del canal $p^{sr}r$. Entonces se puede mostrar que para cada error de configuración alcanzable en el sistema original, hay una configuración que presenta el mismo error que es alcanzable en el sistema emulado y viceversa.
\begin{itemize}
\item[$\Rightarrow$] Consideremos cualquier configuración de error $e$ alcanzable en el sistema original y consideremos cualquier camino que alcance $e$. Luego reemplacemos toda acción $sr!m$ con la secuencia de acciones $sp^{sr}!m, sp^{sr}?m, p^{sr}r!m$ y cada acción $sr?m$ con $p^{sr}?m$, este nuevo camino está presente en el sistema emulado y alcanza el estado donde las configuraciones de buffer y transiciones permitidas son las mismas para el conjunto de máquinas compartidas (es decir, todas las máquinas de sistema emulado menos los repetidores) y los repetidores están en su estado inicial con buffers vacíos.
\item[$\Leftarrow$] Consideremos cualquier error de configuración $e$ alcanzable en el sistema emulado y ocnsideremos cualquier camino que alcance $e$. Ya establecimos que el error no puede ser culpa de los repetidores dado que siempre pueden progresar. Por lo tanto sin importar el estado de los buffers en $e$ hay una configuración $e'$ que presenta el mismo error de comunicación y donde el estado de los buffers de los repetidores está vacío y están en su estado inicial. Luego consideremos que cualquier camino que alcance $e'$ está formado por secuencias de la forma $sp^{sr}!m,\ldots, sp^{sr}?m, \ldots, p^{sr}r!m$ y $p^{sr}r?m$ (donde los puntos suspensivos denota la posibilidad de interleaving de otras acciones). Entonces alcanza con reemplazar cada secuencia $sp^{sr}!m,\ldots_0, sp^{sr}?m, \ldots_1, p^{sr}r!m$ con la secuencia $\ldots_0,\ldots_1, sr!m$ y cada $p^{sr}r?m$ con $sr?m$ para obtener un camino que existe en el sistma original y aclanza una configuración que tiene sus buffers en el mismo estado con las mismas transiciones permitidas.
\end{itemize}
\end{proof}
\end{prop}
%!TEX root = ./main.tex
\cnote{Este texto parece muy introductorio, habla de la motivación para introducir AFCAs y mCFSM pero eso ya pasó en las secciones anteriores. Pasar este texto como introducción del capítulo 3 diciendo para cada cosa dónde lo presentás.}

Como ya hemos mencionado, en SOC los sistemas son concebidos como objetos dinámicos construidos en \emph{run-time} en la medida que su ejecución llega a un estado en el que la intervención de servicios externos se hace necesaria. Es decir un sistema de este tipo utilizará distintos servicios según las necesidades que se manifiesten a lo largo de una ejecución particular. 
 
Una aplicación que se encuentra ejecutando se conecta con los servicios que le son necesarios a través de canales de comunicación por los cuales se envían o reciben mensajes. Estos canales pueden establecer una comunicación entre un número fijo (para cada canal particular) pero no acotado a priori de servicios. En la sección anterior hemos detallado un conjunto de propiedades que garantizan una comunicación sin errores a través de estos canales (i.e. ausencia de deadlock, ausencia de mensajes huérfanos y ausencia de situaciones en las que el receptor no se encuentra a la espera de un mensaje que le fue enviado) y un procedimiento para garantizarlas. Estas condiciones y su procedimiento de análisis parten de la hipótesis de que las CFSMs correspondientes a todos y cada uno de los participantes de la comunicación sobre dicho canal se encuentran disponibles.

Ahora bien, que la aplicación arribe a un estado en el que un servicio se hace necesario sobre un canal particular, no implica que todos los participantes también lo sean en ese mismo instante y por ello, con el objeto de profundizar esta concepción incremental, a demanda, que se tiene sobre los sistemas de software surge, más o menos naturalmente, la idea de poder dotar al \emph{middleware} de la capacidad de realizar un \emph{binding} parcial sobre los canales. A esta práctica la llamaremos \emph{binding incremental}.
 
Esta percepción parcial del \emph{binding} sobre un canal requiere la utilización de un un lenguaje de descripción que soporten dichos mecanismos de composición. Por ejemplo, al componer dos CFSMs puede ocurrir que cada máquina se comunique a través de un canal con un tercer participante, estos dos canales son independientes y, además, una vez que se ha realizado la composición, deben ser percibidos por este tercer participante como canales de comunicación separados. Por lo tanto, para preservar la semántica de la comunicación, es necesario que el CFSM resultante de la composición tenga dos canales con este tercer participante. Este fenómeno será el eje rector de las modificaciones que introduciremos en esta sección. Esta característica no solo será necesario a nivel de CFSM (interfaz de comunicación de un servicio) sino también de los autómatas que caracterizan el cómputo, cuya interfaz de comunicación es expresada a través de una CFSM.

Para modelar este comportamiento introducimos los Autómatas Finitos de Comunicación Asincrónica (AFCA). Estos autómatas tienen tres tipos de transiciones: 
\begin{inparaenum}[1)]
\item internas, que sirven el propósito de representar cómputo realizado por la componente; 
\item de buffer, que representan comunicación asincrónica interna entre elementos de la componente, y que permiten representar la comunicación entre dos participantes luego de una composición; y por último 
\item de entrada / salida, que modelan acciones de comunicación con otras componentes del sistema.
\end{inparaenum} \\

\cnote{Abrir una sección sobre equivalencia del comportamiento de la composición de los AFCAs y el del communicating systems formado por sus interfaces de comunicación.}
Esta sección va a estar enfocada en la definición de composición total de autómatas y la noción de que dado un conjunto $\{ \mathcal{A}_i\}_{i \in I}$ de AFCA weak deterministic sin transiciones de buffer la semántica de la composición es equivalente a la del conjunto de mCFSM proyectadas de los elementos de $\{ \mathcal{A}_i\}_{i \in I}$. 

\begin{figure}[ht]
$$
\xymatrix{   
	\{\mathcal{A}_i\}_{i \in I} \ar[rrr]_{\Pi} \ar[d]_{||} & & & \{C_i\}_{i \in I}  \ar[d]_{semantica}  \\
	  {A} \ar[r]_{\Pi'} \ar[d]_{\Pi} & \mathcal{L}' & = & {\mathcal{L}}  \\
	  \emptyset
}
$$
\caption{La composición de componentes asincrónicas preserva la semántica asincrónica de su interfaz de comunicación.}
\label{fig:preservacion}
\end{figure}

En general, cuando se trata de formalismos que modelan simultáneamente aspectos del comportamiento interno de una componente y elementos de interfaz observable por otras componentes, la proyección de dicha interfaz no resulta bisimilar con el comportamiento real de la componente. A continuación se muestra un ejemplo de este fenómeno recurriendo a la clase de autómatas presentada en la \ref{def:afca}, tomando como comportamiento interno las transiciones internas y de buffer y como comportamiento observable las acciones de comunicación, tanto de entrada como de salida, con otras componentes.

\cnote{Armar una sola figura y usar subfloats para esta figura. El ejemplo no está bien, te mando una foto.}
\begin{figure}[ht]
\begin{tikzpicture}[->, thick]
 \node[state,initial] (q_0)   {$q_0$}; 
 \node[state] (q_1) [below = of q_0 ] {$q_1$};
 \node[state] (q_2) [right = of q_1 ] {$q_2$};
 \node[state] (q_3) [below = of q_1 ] {$q_3$};
 \node[state] (q_4) [below = of q_2 ] {$q_4$};
 \node[state, accepting] (q_5) [below = of q_3 ] {$q_5$};
 \draw[]        
        (q_0) edge[left] node{int1} (q_1)
        (q_0) edge[right] node{int2} (q_2)
        (q_1) edge[left] node{in(sp,a)} (q_3)
        (q_3) edge[left] node{in(sp,b)} (q_5)
        (q_2) edge[right] node{in(sp,b)} (q_4)
        ; 
\end{tikzpicture}
\qquad
\begin{tikzpicture}[->, thick]
 \node[state,initial] (q_0)   {$q_0$}; 
 \node[state] (q_1) [below = of q_0 ] {$q_1$};
  \node[state, accepting] (q_2) [below = of q_1 ] {$q_2$};
 \draw[]        
        (q_0) edge[left] node{out(sp,a)} (q_1)
        (q_1) edge[left] node{out(sp,b)} (q_2)
        ; 
\end{tikzpicture}

\caption{Los AFCA P y S, donde P no es Weak deterministic}
\label{fig:weak determinacy afca}
\end{figure}

\begin{figure}[ht]
    \centering
    
\begin{tikzpicture}[->, thick]
 \node[state,initial] (q_0)   {$q_0$}; 
  \node[state] (q_1) [below = of q_0 ] {$q_1$};
 \node[state] (q_2) [right = of q_1 ] {$q_2$};
 \node[state, accepting] (q_3) [below = of q_1 ] {$q_3$};
 \draw[]        
        (q_0) edge[left] node{sp?a} (q_1)
        (q_0) edge[right] node{sp?b} (q_2)
        (q_1) edge[left] node{sp?b} (q_3)
        ; 
\end{tikzpicture}
\qquad
\begin{tikzpicture}[->, thick]
 \node[state,initial] (q_0)   {$q_0$}; 
 \node[state] (q_1) [below = of q_0 ] {$q_1$};
  \node[state, accepting] (q_2) [below = of q_1 ] {$q_2$};
 \draw[]        
        (q_0) edge[left] node{sp!a} (q_1)
        (q_1) edge[left] node{sp!b} (q_2)
        ; 
\end{tikzpicture}

   \caption{Las mcfsm correspondientes a los autómatas S y P}
    \label{fig:weak determinacy mcfsm}
\end{figure}

\cnote{Explicar acorde al nuevo ejemplo.} En la Fig 3.4\cnote{Usar referencia!} podemos ver al sistema compuesto de los AFCA $S$ y $P$. En la misma podemos observar como $S$ tiene un único camino lineal donde envía los mensajes $a$ y $b$ al autómata $P$. Por su lado $P$ puede tomar dos caminos, ejecutando una de sus dos transiciones internas. Esas transiciones no afectan directamente el comportamiento del autómata $S$. Pero cual de ellas elija ejecutar $P$, sí lo afecta. Al tomar el primer camino el sistema se ejecuta de manera correcta, pero si ejecuta la transición interna $int2$, no puede recibir el mensaje $a$ y por lo tanto el sistema entra en deadlock. El autómata $P$ es un autómata determinístico, pero esa noción no alcanza para que el sistema funciones correctamente. En la figura 3.5 se ve como la mCFSM proyectada preserva los errores comunicacionales.  

Para prevenir este tipo de problemas surgen la noción de determinismo débil o weak determinacy. La noción de determinismo es importante dado que está relacionada con la predecibilidad del sistema. Si realizamos el mismo experimento dos veces sobre un sistema determinista, comenzando ambas veces del estado inicial, esperamos obtener el mismo resultado o comportamiento cada vez. Además, la predecibilidad es algo que frecuentemente requerimos de sistemas. Por ende si elegimos especificar un sistema $S$ dándole un participante abstracto $Spec$ al que $S$ debería ser equivalente, entonces $Spec$ va a ser determinista. La idea de \emph{Weak Determinacy} \cite[Chap.~11, Def.~3]{milner89}\cnote{subir esta cita al primer lugar en que lo mencionás.} es una forma más amplia de determinismo que nos permite contemplar una equivalencia observacional entre participantes. Esta equivalencia basada en la definición de bisimulación tiene en cuenta que para un participante externo, las transiciones $\epsilon$ son invisibles y por lo tanto una serie de estados conectados por transiciones consecutivas de este tipo serían bisimilares. En este caso tomamos las transiciones internas o de comunicación interna como observacionalmente bisimilares para un participante externo.  

%\begin{definition}[Determinacy] 
%Dado un autómata $ \Lambda = \langle Q, \Sigma, \delta, q_0, F \rangle$ decimos que es \emph{determinate} si vale que $(\forall q, q', q'' \in Q)(\nexists t \in \Sigma)(q \xrightarrow{t} q' \land q \xrightarrow{t} q''$.
%\end{definition}

\begin{definition}[Weak Determinacy]
Sea $A= \langle Q, \mathcal{C}, B, \Sigma, \delta, q_0, F \rangle$ un AFCA y $\Sigma' \subseteq \Sigma$ un conjunto de etiquetas a silenciar. Se dice que $A$ es \emph{weak determinate sobre $\Sigma'$} si para todo $q, q', q'' \in Q$ y para todo $s, s' \in \Sigma^*$ tal que $\widehat{s}_{\Sigma'} = \widehat{s'}_{\Sigma'}$, $q \xrightarrow{s} q'$ y $q \xrightarrow{s'} q''$, $q' \sim_{\Sigma'} q''$.
\end{definition}

\begin{prop}
Dado un AFCA $A = \langle Q, \mathcal{C}, B, \Sigma, \delta, q_0, F \rangle$ weak determinate tal que $\Sigma_\mathit{Int} \subseteq \Sigma$ denota el conjunto de transiciones internas, entonces dados dos estados $q, q' \in Q$ tales que $q \xrightarrow{\sigma} q'$ con $\sigma \in \Sigma_{\mathit{Int}}^*$,  $q \sim_{\Sigma_\mathit{Int}} q'$.
\end{prop}
\begin{proof}
Supongamos que $A$ es weak determinate sobre $\Sigma_\mathit{Int}$ y existen $q, q' \in Q$ tales que $q \xrightarrow{\sigma} q'$ con $\sigma \in \Sigma_{\mathit{Int}}^*$ y no vale que $q \sim_{\Sigma_\mathit{Int}} q'$. Si $q \sim_{\Sigma_\mathit{Int}} q'$ no vale, entonces no existen $q'' \in Q$,  $s, s' \in \Sigma^*$ tal que $\widehat{s}_{\Sigma'} = \widehat{s'}_{\Sigma'}$, $q \xrightarrow{s} q'$ y $q \xrightarrow{s'} q''$, lo que es absurdo pues se puede tomar $q'' = q$, $s = \epsilon$ y $s' = s\sigma$ satisfaciendo que $\widehat{s}_{\Sigma_\mathit{Int}} = \widehat{s'}_{\Sigma_\mathit{Int}}$
\end{proof}

\begin{definition}[Equivalencia por comunicación interna]
\cnote{Modificada.}Dado un AFCA $A = \langle Q, \mathcal{C}, B, \Sigma, \delta, q_0, F \rangle$ weak determinate tal que $\Sigma_\mathit{Int} \subseteq \Sigma$ denota el conjunto de transiciones internas, definition $q \Rightarrow q'$ si y solo si existe $\sigma \in \Sigma_\mathit{Int}$ tal que $q \xrightarrow{\sigma} q'$. Luego, $[q]_m$, \emph{la clase de equivalencia por comunicación interna} de $q$, se define como $[q]_m = \{q'\ |\ q \Rightarrow^\bullet q'\}$. 
\end{definition}

Intuitivamente, la clase de equivalencia por comunicación interna de un estado $q$, $[q]_m$ son aquellos estados $q'$ relacionados en cero o más pasos de transiciones internas con $q$.

Dado un AFCA $A = \langle Q, \mathcal{C}, B, \Sigma, \delta, q_0, F \rangle$, $\Sigma = \Sigma_\mathit{Int} \cup \Sigma_\mathit{Buff}$ y $q \in Q$, . Decimos que dos clases de equivalencia están relacionadas $[q]_m \xrightarrow{\sigma} [q']_m$ sii $\exists q_i \in [q]_m, q'_i \in [q']_m$ tales que $q_i \xrightarrow{\sigma} q'_i \in \delta_{\text{Buff}}$.

\begin{definition}[Proyección de comunicación interna]\label{def:pci}\ \\
Dado un AFCA $\mathcal{A} = \langle Q, B, \mathcal{C}, \Sigma, \delta, q_0, F \rangle$, definimos $\Tau(\mathcal{A})$ como la proyección de la semántica de comunicación interna de $\mathcal{A}$. Llamamos $\widehat{\mathcal{A}}$ al sistema de transición etiquetado resultante y decimos que $\Tau (\langle Q, B, \mathcal{C}, \Sigma, q_0, \delta, F \rangle)= \langle \widehat{Q}, B, \widehat{q_0}, \widehat{\Delta} \rangle$ tal que:
\begin{itemize}
    \item $ \widehat{Q} = \{ \langle [q]_m, \overrightarrow{\Omega} \rangle | q \in Q \land \overrightarrow{\Omega} = (\Omega_b)_{b \in B}$ con $\Omega_b \in \mathcal{M}^* \}$ Cada estado del sistema de transición resultante es una configuración de $\mathcal{A}$ para una clase de equivalencia por comunicación interna,
    \item $B$ es el conjunto de buffers del LTS, que son los mismos de comunicación interna de $\mathcal{A}$,
    \item $\widehat{q_0}= \langle q_0, \Omega \rangle$ es el estado inicial de  donde $q_0$ es el estado inicial de $\mathcal{A}$ y $\Omega = \langle [], \ldots, [] \rangle$ es el conjunto de buffers en su estado inicial sin mensajes, y
    \item $\widehat{\Delta}= \{\langle \langle q_i, \Omega \rangle, \widehat{\sigma},\langle q_j, \Omega' \rangle \rangle | q_i, q_j \in Q, \widehat{\sigma} \in \{B \times \{\gg, \ll\} \times \mathcal{M}\} \land \exists \sigma \in \Sigma$ tal que  $[q_i]_m \xrightarrow{\sigma} [q_j]_m \}$. Definimos $\widehat{\Delta}$ como la relación de transición de $\widehat{\mathcal{A}}$. Un estado de $\widehat{q_j}$ es alcanzable desde otro $\widehat{q_i}$ si la clase de equivalencia $[q_j]_m$ es alcanzable desde $[q_i]_\mathcal{M}$.
\end{itemize}
\end{definition}

Antes de pasar a la demostración de equivalencia entre la semánticas de un conjunto compatible de AFCAs

\begin{lemma}[Equivalencia de estados discretos] Sea $\widehat{\mathcal{A}}=\langle \widehat{Q}, B, \widehat{q_0}, \widehat{\delta} \rangle$, la proyección de comunicación interna del AFCA $\mathcal{A}$ que surge de la composición del conjunto de autómatas  $\{\mathcal{A}\}_{1 \leq i \leq n}$  y $M= \langle [Q], \mathcal{C}, [q_0], [\delta] \rangle$, el LTS que representa la semántica del CS correspondiente a \cnote{Poner el conjunto de las proyecciones de los autómatas.}. Para todo $\langle [\overrightarrow{q}]_m, \overrightarrow{\Omega} \rangle \in \widehat{Q}$, existe $\langle \overrightarrow{p}, \overrightarrow{\Omega_{\mathcal{C}}} \rangle \in P$ tal que $q_i \in \overrightarrow{q} \iff q_i \in \overrightarrow{p}$ y viceversa. 
\end{lemma}
\begin{proof} 
\begin{itemize}
    \item[$\implies$] Dado $\widehat{q} \in \widehat{Q}$, un estado de la proyección $\widehat{\mathcal{A}}$. Por definición de $\widehat{\mathcal{A}}$, $\widehat{q}= \langle [\overrightarrow{q}]_m, \overrightarrow{\Omega} \rangle$. Por definición de equivalencia por comunicación interna $[\overrightarrow{q}]_m = \left\{\overrightarrow{q'}\ |\ \overrightarrow{q} \xrightarrow{\equiv}\overrightarrow{q'}\right\}$\cnote{qué es este ``equiv'', cambiar usando la notación e imagino que usando $\sigma$.}  con $\overrightarrow{q},\overrightarrow{q'} \in Q$ y $\sigma \in \Sigma_\mathit{Int}^*$ o sea son los estados alcanzables por transiciones internas desde $\overrightarrow{q}$. Por lo tanto $\langle [\overrightarrow{q}]_m, \overrightarrow{\Omega} \rangle \in \widehat{Q} \iff \overrightarrow{q} \in \Q$. Por definición de composición de AFCA $\overrightarrow{q} = \langle q_1, \ldots, q_i, \ldots, q_n \rangle$ donde $q_i \in Q_i$ es un estado del AFCA $\mathcal{A}_i$. Por definición de proyección de comunicación interna $\overrightarrow{q}$ es un estado alcanzable por una transición de comunicación interna del AFCA $\mathcal{A}$. Como el conjunto de AFCA $\{\mathcal{A}\}_{1 \leq i \leq n}$ no posee comunicación interna, la comunicación interna de $\mathcal{A}$ es la comunicación externa de sus componentes internalizada. Por lo tanto $q_i$ es un estado alcanzable por una transición de comunicación externa de $\mathcal{A}_i$. Entonces$\overrightarrow{q} \in Q \iff [q_i] \in Q_i$ y por definición de interfaz de comunicación de AFCA $[q_i] \in [Q_i]$ que es el conjunto de estados de la mCFSM $M_i$. Por lo tanto existe $\overrightarrow{p} = \langle [q_1], \ldots, [q_i], \ldots, [q_n] \rangle$ tal que $\langle \overrightarrow{p}, \overrightarrow{\Omega_\mathcal{C}} \rangle \in P$.
    
    \item[$\impliedby$] Sea $p \in P$, un estado del LTS $M$ que representa la semántica el CS. Por definición de CS vale $p = \langle \langle [q_1],\ldots,[q_i], \ldots, [\q_n] \rangle, \overrightarrow{\Omega} \rangle \in P \iff [q_i] \in [Q_i]$, donde $[Q_i]$ es el conjunto de estados de la mCFSM $M_i$. Sea $\mathcal{A}_i$ el AFCA tal que $M_i$ es su interfaz de comunicación. Por definición de interfaz de comunicación existe un estado alcanzable por una transición de comunicación externa $q_i \in Q_i$. Por definición de composición $q_i \in Q_i \iff \langle q_1, q_2, \ldots, q_i, \ldots, q_n \rangle \in Q$ y la transición de comunicación externa es internalizada. Por lo tanto $\overrightarrow{q}$ es alcanzable por una transición de comunicación interna y tiene su clase de comunicación interna $[\overrightarrow{q}]_m$. Entonces existe $\widehat{q} = \langle [\overrightarrow{q}]_m, \overrightarrow{\Omega} \rangle \in \widehat{Q}$
\end{itemize}
\end{proof}

\begin{lemma}[Equivalencia entre buffers y canales]\label{lemma2} Sean $\widehat{\mathcal{A}}=\langle \widehat{Q}, B, \widehat{q_0}, \widehat{\delta} \rangle$, la proyección de comunicación interna del AFCA $\mathcal{A}$ que surge de la composición del conjunto de autómatas  $\{\mathcal{A}\}_{1 \leq i \leq n}$  y $M= \langle P, \mathcal{C}, p_0, [\delta] \rangle$, el LTS que representa la semántica del CS correspondiente a las mCFSM del conjunto. Dados $\langle [\overrightarrow{q}]_m, \overrightarrow{\Omega} \rangle \in \widehat{Q}$ y $\langle \overrightarrow{p}, \overrightarrow{\Omega_{\mathcal{C}}} \rangle \in P$. Queremos ver que el conjunto de buffers $B$ de $\mathcal{A}$ es el conjunto de canales internalizados de $\mathcal{C}$ de $M$. Es decir que para todo canal $\omega$ vale que $\omega \in \mathcal{C} \iff \omega \in B$.
\end{lemma}
\begin{proof} Sea $\langle \overrightarrow{p}, \overrightarrow{\Omega_{\mathcal{C}}} \rangle$ un estado del LTS $M$. Por definición de CS $\overrightarrow{\Omega_{\mathcal{C}}}$ es el conjunto de canales $\mathcal{C}$ con su contenido en el estado discreto $\overrightarrow{p}$. Estos canales son los canales de las mCFSM $\{M_i\}_{1 \leq i \leq n}$ que componen el CS. Entonces vale $(\omega \in \overrightarrow{\Omega_{\mathcal{C}}}) \iff (\omega \in \mathcal{C}) \iff (\exists M_i \in \{M_i\}_{1 \leq i \leq n} \land  \omega \in \mathcal{C}_i)$. Sea $\mathcal{A}_i$ el AFCA tal que $M_i$ es su interfaz de comunicación. Por definición de interfaz de comunicación de AFCA $\omega \in \mathcal{C}_i \iff \omega \in {\mathcal{C}}_{\mathcal{A}_i}$. Por definición de composición de AFCA, la comunicación externa se internaliza. Entonces como $\mathcal{A}_i$ no tienen buffers ni comunicación interna, vale $\omega \in {\mathcal{C}}_{\mathcal{A}_i} \iff \omega \in B$ donde $B$ es el conjunto de buffers de comunicación interna de $\mathcal{A}$. Luego por definición de proyección de comunicación interna el conjunto de buffers de $\mathcal{A}$ es el mismo que en $\widehat{\mathcal{A}}$. Por lo tanto $\omega \in B \iff \exists \langle [\overrightarrow{q}]_m, \overrightarrow{\Omega} \rangle \in \widehat{Q}$ tal que $\omega \in \overrightarrow{\Omega}$.
\end{proof}


\begin{theorem}[Equivalencia de las semánticas de las mFCSMs y los AFCAs] Sea $\{\mathcal{A}\}_{1 \leq i \leq n}$ un conjunto de weak deterministic AFCA sin transiciones de comunicación interna, tal que su composición $\mathcal{A} = \langle Q, B, \mathcal{C},\Sigma, \delta, q_0, F \rangle$, también lo es. Llamamos $\widehat{\mathcal{A}}=\langle \widehat{Q}, B, \widehat{q_0}, \widehat{\delta} \rangle$ al sistema de transición etiquetado (LTS) que resulta de la proyección $\Tau(\mathcal{A})$. Sea el conjunto $\{M_i\}_{1 \leq i \leq n} \langle [Q_i], \mathcal{C}, [{q_0}_i], [\delta_i] \rangle$ de conformado por las mCFSM correspondientes al conjunto AFCA. Llamamos $M = \langle P, \mathcal{C}, {p_0}, [\delta] \rangle$ al LTS que representa la semántica del CS. Luego $\widehat{\mathcal{A}}$ y $M$ son bisimilares. 
\end{theorem}

\begin{proof}
Queremos probar que la relación $R \subseteq \widehat{Q} \times P$, definida de la siguiente manera $R = \{ \langle \langle [\overrightarrow{q}]_m, \overrightarrow{\Omega} \rangle, \langle \overrightarrow{p}, \overrightarrow{\Omega_{\mathcal{C}}} \rangle \rangle \ | \ [\overrightarrow{q}]_mR'\overrightarrow{p} \land  \overrightarrow{\Omega}=\overrightarrow{\Omega_{\mathcal{C}}} \}$ es una bisimulación, con $R'=\{\langle [\overrightarrow{q}]_m, \overrightarrow{p} \rangle \ | \ q_i \in \overrightarrow{q} \iff q_i \in \overrightarrow{p} \}$. Donde $P$ es el conjunto de estados de $M$ y para todo $\langle [\overrightarrow{q}]_m, \overrightarrow{\Omega} \rangle \in \widehat{Q}$, $\langle \overrightarrow{p}, \overrightarrow{\Omega_{\mathcal{C}}} \rangle \in P$, $\langle \langle [\overrightarrow{q}]_m, \overrightarrow{\Omega} \rangle, \langle \overrightarrow{p}, \overrightarrow{\Omega_{\mathcal{C}}} \rangle \rangle \in R$ si y solo si $[\overrightarrow{q}]_mR'\overrightarrow{p}$ y $\overrightarrow{\Omega}=\overrightarrow{\Omega_{\mathcal{C}}}$. \\

Es decir que para cada par de elementos $\widehat{q} \in \widehat{Q}$, $p \in P$, vale $\langle \widehat{q}, p \rangle \in R$ y $\forall \widehat{\sigma} \in \widehat{\Sigma} | \widehat{q} \xrightarrow{\widehat{\sigma}} \widehat{q}' \implies \exists p' \in P$ tal que $p \xrightarrow{\sigma_M} p' \land \langle \widehat{q}', p', \rangle \in R$, y simétricamente $\forall \sigma_M \in \Sigma_M | p \xrightarrow{\sigma_M} p \implies \exists \widehat{q}' \in \widehat{Q}, \widehat{\sigma} \in \widehat{\Sigma}$ tal que $\widehat{q} \xrightarrow{\widehat{\sigma}} \widehat{q}' \land \langle \widehat{q}', p'\rangle \in R$. Donde $\widehat{\Sigma} \subseteq B \times \{\gg,\ll\} \times \mathcal{M}$, $\Sigma_M \subseteq \mathcal{A} \times \{!,?\} \times \mathcal{M}$ y el conjunto de buffers $B$ son los canales internalizados del conjunto $\mathcal{C}$. 

Tomamos inicialmente como etiquetas $\widehat{\sigma} = (a_ka_j)_l \gg m$ y $\sigma_M = (a_ka_j)_l?m$, 

Primero demostraremos la condición $\forall (a_ka_j)_l \gg m \in \widehat{\Sigma} | \widehat{q} \xrightarrow{(a_ka_j)_l \gg m} \widehat{q}' \implies \exists p' \in P$ tal que $p \xrightarrow{(a_ka_j)_l?m} p' \land \langle \widehat{q}', p', \rangle \in R$. 
 
\begin{itemize}
    
    \item Dado el par $\langle [\overrightarrow{q}]_m, \overrightarrow{\Omega} \rangle \in \widehat{Q}, \langle \overrightarrow{p}, \overrightarrow{\Omega_\mathcal{C}} \rangle \in P$ tal que $\langle \langle [\overrightarrow{q}]_m, \overrightarrow{\Omega} \rangle, \langle \overrightarrow{p}, \overrightarrow{\Omega_\mathcal{C}} \rangle \rangle \in R$. Sean $(a_ka_j)_l \gg m \in \widehat{\Sigma}$ y $\langle [\overrightarrow{q}']_m, \overrightarrow{\Omega}' \rangle \in \widehat{Q}$ tales que $\langle [\overrightarrow{q}]_m, \overrightarrow{\Omega} \rangle \xrightarrow{(a_ka_j)_l \gg m} \langle [\overrightarrow{q}']_m, \overrightarrow{\Omega}' \rangle \in \widehat{\Delta}$.  Esto vale si y solo si la clase de equivalencia por comunicación interna $[\overrightarrow{q}']_m$ es alcanzable desde $[\overrightarrow{q}]_m$ vía $(a_ka_j)_l \gg m$. Es decir $[\overrightarrow{q}]_m \xrightarrow{(a_ka_j)_l \gg m} [\overrightarrow{q}']_m$. Entonces $\overrightarrow{\Omega} = [(a_1a_2), \ldots, (a_ka_j)_l \cdot m \ldots (a_na_m)]$ y $\overrightarrow{\Omega}' = [(a_1a_2), \ldots, (a_ka_j)_l \ldots (a_na_m)] \rangle$.
    
    \item Sean $\overrightarrow{q_{\sigma}} \in [\overrightarrow{q}]_m, \overrightarrow{q_{\sigma}}' \in [\overrightarrow{q}']_m$, por definición de proyección de comunicación vale $\overrightarrow{q_{\sigma}} \xrightarrow{(a_ka_j)_l \gg m} \overrightarrow{q_{\sigma}}' \in \delta$. Por definición de composición sabemos que vale $\overrightarrow{q_{\sigma}} =\langle q_0, q_1, \ldots, q_j, \ldots,q_n \rangle$ y $\overrightarrow{q_{\sigma}}' =\langle q_0, q_1, \ldots, q'_j, \ldots,q_n \rangle$ donde $q_j,q'_j$ son estados de $\mathcal{A}_j$, componente de $\mathcal{A}$. Por definición de AFCA, $(a_ka_j)_l \gg m$ es una operación de desencolamiento del mensaje $m$ sobre el buffer $(a_ka_j)_l$. Como los AFCA del conjunto no tienen acciones internas esto corresponde a una recepción de un mensaje por parte de $\mathcal{A}_j$ vía el canal $(a_ka_j)_l$. Entonces podemos afirmar que se cumple $\overrightarrow{q_{\sigma}} \xrightarrow{(a_ka_j)_l \gg m} \overrightarrow{q_{\sigma}}'\in \delta \iff q_j,q'_j \in Q_j \land q_j \xrightarrow{(a_ka_j)_l?m} q'_j \in \delta_j$.
    
    \item Por la interfaz de comunicación de AFCA podemos obtener $M_j$, la mCFSM correspondiente a $\mathcal{A}_j$, y sabemos que $q_j \xrightarrow{(a_ka_j)_l?m} q'_j \in \delta_j \iff$ se cumple que $[q_j], [q'_j] \in [Q_j] \land [q_j] \xrightarrow{(a_ka_j)_l?m} [q'_j] \in [\delta_j]$. 
    
    \item Entonces, por definición del LTS M, existen $p'=\langle \overrightarrow{p'}, \overrightarrow{\Omega_\mathcal{C}} \rangle \in P, (a_ka_j)_l?m \in \Sigma_M$ tal que $\overrightarrow{p'}=\langle [q_1], [q_2], \ldots, [q'_j], \ldots [q_n] \rangle$, $\overrightarrow{\Omega'_\mathcal{C}}=[\ldots, (a_ka_j)_l, \ldots]$, y $p \xrightarrow{(a_ka_j)_l?m} p'$. Donde $\overrightarrow{\Omega_\mathcal{C}}$ y $\overrightarrow{\Omega_\mathcal{C}}'$ son los canales de comunicación externa para los estados $\overrightarrow{p}$ y $\overrightarrow{p'}$, y vale $\overrightarrow{\Omega_\mathcal{C}} = [(a_1a_2), \ldots, (a_ka_j)_l \cdot m \ldots (a_na_m)]$ y $\overrightarrow{\Omega_\mathcal{C}}' = [(a_1a_2), \ldots, (a_ka_j)_l \ldots (a_na_m)] \rangle$. Por lo tanto vale $\langle \widehat{q'}, p' \rangle \in R$.
    
    \item Por lo tanto vale que $\forall (a_ka_j)_l \gg m \in \widehat{\Sigma} | \widehat{q} \xrightarrow{(a_ka_j)_l \gg m} \widehat{q}' \implies \exists p' \in P$ tal que $p' \xrightarrow{(a_ka_j)_l?m} p' \land \langle \widehat{q}', p', \rangle \in R$ donde $\widehat{\Sigma} \subseteq B \times \{\gg,\ll\} \times \mathcal{M}$
  \end{itemize}
  
  Demostraremos ahora la condición simétrica $\forall (a_ka_j)_l?m \in \Sigma_M | p \xrightarrow{(a_ka_j)_l?m} p' \implies \exists \widehat{q}' \in \widehat{Q}, (a_ka_j)_l \gg m \in \widehat{\Sigma}$ tal que $\widehat{q} \xrightarrow{(a_ka_j)_l \gg m} \widehat{q}' \land \langle \widehat{q}', p'\rangle \in R$ 
  \begin{itemize}

     \item Dado el par $\langle q, p \rangle = [\overrightarrow{q}]_m, \overrightarrow{\Omega} \rangle \in \widehat{Q}, \langle \overrightarrow{p}, \overrightarrow{\Omega_\mathcal{C}} \rangle \in P$ tal que $\langle \langle [\overrightarrow{q}]_m, \overrightarrow{\Omega} \rangle, \langle \overrightarrow{p}, \overrightarrow{\Omega_\mathcal{C}} \rangle \rangle \in R$. Sean $(a_ka_j)_l?m$ y $p' = \langle \overrightarrow{p}', \overrightarrow{\Omega_\mathcal{C}}' \rangle \in P$ tales que $\langle \overrightarrow{p}, \overrightarrow{\Omega_\mathcal{C}} \rangle \xrightarrow{(a_ka_j)_l?m} \langle \overrightarrow{p}', \overrightarrow{\Omega_\mathcal{C}}' \rangle \in \delta_P$. Por definición del LTS $M$, $(a_ka_j)_l?m$ es una operación sobre el canal $(a_ka_j)_t$ de la mCFSM $M_j$, componente del CS. Como el LTS $M$ representa la semántica del CS compuesto por el conjunto de mCFSMs, vale que $\overrightarrow{p}=\langle [q_1], [q_2], \ldots [q_j], \ldots [q_n] \rangle$, $\overrightarrow{p}'= \langle [q_1], [q_2], \ldots, [q'_j], \ldots [q_n] \rangle$, donde $[q_j]$ y $[q'_j]$ son estados de la mCFSM $M_j$, y vale $[q_j] \xrightarrow{(a_ka_j)_l?m} [q'_j] \in [\delta]_j$
    
    \item Por definición de intefaz de comunicación de AFCA sabemos que $M_J$ es la mCFSM correspondiente a $\mathcal{A}_j$, componente de $\mathcal{A}$. Por lema 1, como $[q_j] \xrightarrow{(a_ka_j)_l?m} [q'_j] \in [\delta_j]$ entonces $q_j \xrightarrow{(a_ka_j)_l?m} q'_j \in \delta_j$.
    
    \item Por definición de composición $q_j$ y $q'_j$ son componentes de dos estados compuestos que llamamos $\overrightarrow{q}_{\sigma}$ y $\overrightarrow{q}_{\sigma}'$, respectivamente. De modo tal que $\overrightarrow{q}_{\sigma} \xrightarrow{(a_ka_j)_l \gg m} \overrightarrow{q}_{\sigma}' \in \delta$. Donde $(a_ka_j)_l$ es el canal internalizado en forma de buffer.
    
    \item Por definición de equivalencia por comunicación interna como $\overrightarrow{q}_{\sigma} \xrightarrow{(a_ka_j)_l \gg m} \overrightarrow{q}_{\sigma}' \in \delta$ entonces existen las clases de equivalencia $[\overrightarrow{q}]_m$, $[\overrightarrow{q}']_m$, tales que $\overrightarrow{q}_{\sigma} \in [\overrightarrow{q}]_m \land \overrightarrow{q}_{\sigma}' \in [\overrightarrow{q}']_m$ y vale $[\overrightarrow{q}]_m \xrightarrow{(a_ka_j)_l \gg m} [\overrightarrow{q}']_m$. Por proyección de comunicación entonces existen $\widehat{q}' = \langle [\overrightarrow{q}']_m, \overrightarrow{\Omega}' \rangle \in \widehat{Q}, (a_ka_j)_l \gg m \in \widehat{\Sigma}$, tales que $\langle [\overrightarrow{q}]_m, \overrightarrow{\Omega} \rangle \xrightarrow{(a_ka_j)_l \gg m} \langle [\overrightarrow{q}']_m, \overrightarrow{\Omega}' \rangle$. Donde, $\overrightarrow{\Omega}$, y $\overrightarrow{\Omega}'$ son los buffers de $\mathcal{A}$ para los estados $\overrightarrow{q}_{\sigma}$ y  $\overrightarrow{q}_{\sigma}'$, respectivamente, y valen $\overrightarrow{\Omega} = [(a_1a_2), \ldots, (a_ka_j)_l \cdot m \ldots (a_na_m)]$ y $\overrightarrow{\Omega}' = [(a_1a_2), \ldots, (a_ka_j)_l \ldots (a_na_m)] \rangle$. Por lo tanto $\langle \widehat{q}', p' \rangle \in R$ 
    
    % \item Por lo tanto vale en forma simétrica que $\forall (a_ka_j)_l?m \in \Sigma_M | p \xrightarrow{(a_ka_j)_l?m} p \implies \exists \widehat{q}' \in \widehat{Q}, (a_ka_j)_l \gg m \in \widehat{\Sigma}$ tal que $\widehat{q} \xrightarrow{(a_ka_j)_l \gg m} \widehat{q}' \land \langle \widehat{q}', p'\rangle \in R$ donde $\Sigma_M \subseteq \mathcal{C} \times \{!,?\} \times \mathcal{M}$
    
\end{itemize}
\end{proof}


\begin{figure}[h]
\begin{example}[Transformación de Autómatas Finitos de Comunicación Asincrónica]
\label{ex:Transformación}
Consideremos los siguientes AFCA 
\begin{center}
\begin{tikzpicture}[->, thick]
 \node[state,initial] (q_0)   {$q_0$}; 
 \node[state] (q_1) [right= 1.5cm of q_0 ] {$q_1$};
 \node[state] (q_2) [below= of q_0 ] {$q_2$};
 \node[state] (q_3) [right= 1.5cm of q_2 ] {$q_3$};
 \node[state, accepting] (q_4) [below= of q_3 ] {$q_4$};
	
 \draw[]        
        (q_0) edge[above] node{out(pr,a)} (q_1)
        (q_0) edge[left] node{in(sp,b)} (q_2)
        (q_1) edge[right] node{in(sp,b)} (q_3)
        (q_2) edge[above] node{out(pr,a)} (q_3)
        (q_3) edge[right] node{$int_p$} (q_4)
        ; 
\end{tikzpicture} 
\qquad
\begin{tikzpicture}[->, thick]
 \node[state,initial] (q_0)   {$q_0$}; 
 \node[state] (q_1) [below= of q_0 ] {$q_1$};
 \node[state, accepting] (q_2) [below= of q_1 ] {$q_2$};
 \draw[]        
        
        (q_0) edge[right] node{$int_r$} (q_1)
        (q_1) edge[right] node{in(pr,a)} (q_2)
        ;
\end{tikzpicture} 
\qquad
\begin{tikzpicture}[->, thick]
 \node[state,initial] (q_0)   {$q_0$}; 
 \node[state] (q_1) [below= of q_0 ] {$q_1$};

 \draw[]        
        
        (q_0) edge[right] node{out(sp,b)} (q_1)
        
        ;
\end{tikzpicture} 

\end{center}

El autómata compuesto resultante es el AFCA $\mathcal{A}$  \\  
\begin{tikzpicture}[->, thick]
 \node[state,initial] (q_0)   {$q_{000}$}; 
 \node[state] (q_1) [right= 1.5cm of q_0 ] {$q_{100}$};
 \node[state] (q_2) [right = 1.5 of q_1 ] {$q_{110}$};
 \node[state] (q_3) [right = 1.5 of q_2 ] {$q_{120}$};
  \node[state] (q_4) [right = 1.5 of q_3 ] {$q_{121}$};
  \node[state] (q_5) [below=  of q_0] {$q_{001}$};
  \node[state] (q_6) [right = 1.5cm of q_5] {$q_{201}$};
  \node[state] (q_7) [right = 1.5 of q_6 ] {$q_{301}$};
  \node[state] (q_8) [right = 1.5 of q_7 ] {$q_{311}$};
  \node[state] (q_9) [right = 1.5 of q_8 ] {$q_{321}$};
 \node[state,accepting] (q_10) [right= 1.5cm of q_9 ] {$q_{421}$};
 \draw[]        
        (q_0) edge[above] node{$PR \ll a$} (q_1)
        (q_1) edge[above] node{$int_r$} (q_2)
		(q_2) edge[above] node{$PR \gg a$} (q_3)
        (q_3) edge[above] node{$SP \ll b$} (q_4) 
        (q_4) edge[right] node{$SP \gg b$} (q_9)
        (q_0) edge[left] node{$SP \ll b$} (q_5)
        (q_5) edge[below] node{$SP \gg b$} (q_6)
        (q_6) edge[below] node{$PR \ll a$} (q_7)
        (q_7) edge[below] node{$int_r$} (q_8)
        (q_8) edge[below] node{$PR \gg a$} (q_9)
        (q_9) edge[below] node{$int_p$} (q_10)
        ;
\end{tikzpicture}
    \caption{Transformación}
    \label{fig:transformacion}
\end{example}
\end{figure}


\begin{figure}[h]
\begin{example}[CS de transformación]
\label{ex:cmst} El CS correspondiente al conjunto de autómatas es el siguiente
\centering

\begin{tikzpicture}[->, thick]
 \node[state,initial] (q_0)   {$q_0$}; 
 \node[state] (q_1) [right= of q_0 ] {$q_1$};
  \node[state] (q_2) [below= of q_0 ] {$q_2$};
 \node[state] (q_3) [right= of q_2 ] {$q_3$};

 \draw[]        
        (q_0) edge[above] node{pr!a} (q_1)
        (q_0) edge[right] node{sp?b} (q_2)
        (q_1) edge[right] node{sp?b} (q_3)
        (q_2) edge[above] node{pr!a} (q_3)
        ; 
\end{tikzpicture} 
\qquad
\begin{tikzpicture}[->, thick]
 \node[state,initial] (q_0)   {$q_0$}; 
 \node[state] (q_1) [below= of q_0 ] {$q_2$};

 \draw[]        
        
        (q_0) edge[right] node{pr?a} (q_1)
        
        ;
\end{tikzpicture} 
\qquad
\begin{tikzpicture}[->, thick]
 \node[state,initial] (q_0)   {$q_0$}; 
 \node[state] (q_1) [below= of q_0 ] {$q_1$};

 \draw[]        
        
        (q_0) edge[right] node{sp!b} (q_1)
        
        ;
\end{tikzpicture} 

La semántica del CS es el sistema de transición etiquetado $\mathcal{L}$
\\
\begin{tikzpicture}[->, thick][scale=0.20]
 \node[state,initial] (q_0)   {$q_{000},\epsilon,\epsilon$}; 
 \node[state] (q_1) [right= 3cm of q_0 ] {$q_{100},a, \epsilon$};
 \node[state] (q_2) [right = 3cm of q_1 ] {$q_{120},\epsilon, \epsilon$};
 \node[state] (q_3) [below = of q_2 ] {$q_{121}, \epsilon,b$};
 \node[state,accepting] (q_4) [below=  of q_3] {$q_{321}, \epsilon,\epsilon$};
 \node[state] (q_5) [below = of q_0] {$q_{001}, \epsilon, b$};
 \node[state] (q_6) [below = of q_5 ] {$q_{201}, \epsilon, \epsilon$};
 \node[state] (q_7) [right= 3cm of q_6 ] {$q_{301}, a, \epsilon$};
 \draw[]        
        (q_0) edge[above] node{$\langle q_{0p},PR!a,q_{1p} \rangle$} (q_1)
        (q_1) edge[above] node{$\langle q_{0r},PR?a,q_{1r} \rangle$} (q_2)
        (q_2) edge[right] node{$\langle q_{0s},SP!b,q_{1s} \rangle$} (q_3)
        (q_3) edge[right] node{$\langle q_{1p},SP?b,q_{3p} \rangle$} (q_4)
        (q_0) edge[left] node{$\langle q_{0s},SP!b,q_{1s} \rangle$} (q_5)
        (q_5) edge[left] node{$\langle q_{0p},SP?b,q_{2p} \rangle$} (q_6)
        (q_6) edge[below] node{$\langle q_{2p},PR!a,q_{3p} \rangle$} (q_7)
        (q_7) edge[below] node{$\langle q_{0r},PR?a,q_{1r} \rangle$} (q_4)
        ;
\end{tikzpicture}
    \caption{Transformación}
    \label{fig:semántica mcfsm}
\end{example}
\end{figure}

\begin{figure}[ht]
    \begin{centering}
      \begin{tikzpicture}[->, thick]
         \node[state,initial] (q_0) {$[q_{000}],\epsilon, \epsilon$}; 
         \node[state] (q_1) [right= 3.5cm of q_0 ] {$[q_{100}], a, \epsilon$};
         \node[state] (q_2) [right = 3.5cm of q_1 ] {$[q_{120}],\epsilon,\epsilon$};
         \node[state] (q_3) [below= of q_2 ] {$[q_{121}],\epsilon, b$};
         \node[state] (q_4) [below=  of q_0] {$[q_{001}],\epsilon, b$};
         \node[state] (q_5) [below = of q_4] {$[q_{201}],\epsilon, \epsilon$};
         \node[state] (q_6) [right = 3.5cm of q_5 ] {$[q_{301}], a, \epsilon$};
         \node[state,accepting] (q_7) [below = of q_3 ] {$[q_{321}],\epsilon, \epsilon$};
    \draw[]        
         (q_0) edge[above] node{$\langle q_{000}, PR \ll a, q_{100} \rangle$} (q_1)
    	 (q_1) edge[above] node{$\langle q_{110}, PR \gg a, q_{120} \rangle$} (q_2)
         (q_2) edge[right] node{$\langle q_{120}, SP \ll b, q_{121} \rangle$} (q_3) 
         (q_3) edge[right] node{$\langle q_{121}, SP \gg b, q_{321} \rangle$} (q_7)
         (q_0) edge[left] node{$\langle q_{000}, SP \ll b, q_{001} \rangle$} (q_4)
         (q_4) edge[left] node{$\langle q_{001}, SP \gg b, q_{201} \rangle$} (q_5)
         (q_5) edge[below] node{$\langle q_{201}, SP \gg b, q_{301} \rangle$} (q_6)
        (q_6) edge[below] node{$\langle q_{311}, PR \ll a, q_{321} \rangle$} (q_7)
         ;
        \end{tikzpicture}
    \caption{Caption}
    \label{fig: transformación loca}
 \end{centering}
\end{figure}

\newpage



%!TEX root = ./main.tex
\newpage
\section{Conclusiones y Trabajo futuro} 

En esta sección resumimos las contribuciones de esta tesis junto con conclusiones y posibles futuras líneas de investigación.


\subsection{Conclusiones}
En esta tesis exploramos y contribuimos a construir un modelo formal para software orientado a servicios. Partimos de la hipótesis que esto es un paso necesario para lograr la implementación de la infraestructura necesaria para llevar a cabo el paradigma de SOC. Existiendo modelos previos de autómatas como las CFSM, y los IO Automata quisimos hacer un modelo nuevo que tenga en cuenta no solo la comunicación asincrónica entre componentes sino también acciones internas, incluyendo comunicac

\subsection{Trabajo futuro: Composición parcial vs composición total}

En esta sección definimos los Autómatas Finitos de Comunicación Asincrónica para modelar la composición parcial de CFSMs. Ahora necesitamos demostrar que esta composición parcial es equivalente a una composición total.

Como sabemos que todas las CFSMs a componer se encuentran en $\mathcal{P}$ podemos decir que conocemos a priori todos los componentes de la composición final. Dado un conjunto finito $\mathcal{P}$ de CFSMs denominamos $ p_1, p_2, p_3, \ldots, p_n$. 

Si componemos $p_1$ y $p_2$ nos quedaría el conjunto $\mathcal{P}_1 =\{ p_{12} \}, p_3, \ldots, p_n \} $, podemos hacer un paso siguiente componiendo $p_{n-1}$ y $p_n$. De esto obtenemos $\mathcal{P_2}_2= \{ p_{12}, p_3, \ldots, p_n-1 \} $. Podemos continuar este proceso hasta llegar a tener un único autómata. %Esto de acá capaz vuela cuando lo dibuje

Decimos que en cada paso de la composición parcial tengo una función suryectiva, pero no inyectiva que garantiza que cada punto del codominio es la composición de al menos dos puntos de la preimagen.

Llamemos TM1 y TM2 a los autómatas resultantes de la composición de todos los elementos de $\mathcal{P}$ por sucesión de composiciones parciales y por composición total respectivamente. Queremos ver que $ q \in Q_{TM1} \iff q \in Q_{TM2} $ y en ambos casos está entre los estados alcanzables. 

Queremos demostrar que
\begin{enumerate}
\item Las configuraciones alcanzables entre tm1 y tm2 son las mismas
\item tm1 es bisimilar a tm2, es decir que para cada configuración las acciones realizables son las mismas
\end{enumerate}




\bibliographystyle{alpha}
\bibliography{bibdatabase}


\end{document}